\chapter{Conformal Field Theory}

\section{What Is Conformal Field Theory?}

The question is what a conformal field theory is, and why such theories
are useful.  The point of view throughout is that a conformal field
theory gives recipes for local, UV-complete quantum field theories.

A local QFT is a quantum-mechanical system equipped with:
\begin{itemize}
  \item a Hilbert space \(\mathcal H\) of physical states;
  \item Poincare symmetry \(U(\Lambda,a)\), generated by
  \[
    \widehat P_\mu,\qquad \widehat J_{\mu\nu},
  \]
  where \(\widehat P_\mu\) is the energy-momentum generator and
  \(\widehat J_{\mu\nu}\) contains angular momenta and boosts;
  \item a Poincare-invariant vacuum \(\ket{\Omega}\);
  \item local field operators \(\widehat\Phi(x)\) obeying Poincare
  covariance and microcausality.
\end{itemize}
In particular, there is a local stress-energy tensor operator
\(T_{\mu\nu}\), with
\[
  \partial_\mu T^{\mu\nu}=0
  \qquad\text{(conservation)}
\]
and
\[
  T^{\mu\nu}=T^{\nu\mu}
  \qquad\text{(Lorentz symmetry)}.
\]
The generators are obtained from \(T^{\mu\nu}\) by
\[
  \widehat P^\mu=\int \dd^{D-1}\vec x\,T^{0\mu},
\]
\[
  \widehat J^{\mu\nu}
  =
  \int \dd^{D-1}\vec x\,
  \left(x^\mu T^{0\nu}-x^\nu T^{0\mu}\right).
\]

Generally \(T_{\mu\nu}(x)\) transforms as a symmetric tensor under
Lorentz symmetry:
\[
  U(\Lambda)T_{\mu\nu}(x)U(\Lambda)^{-1}
  =
  \Lambda^\rho{}_\mu\Lambda^\sigma{}_\nu
  T_{\rho\sigma}(\Lambda x).
\]
The trace \(T^\mu{}_\mu(x)\) is a scalar operator, and can be expanded
on a basis of scalar field operators \(\mathcal O_I(x)\):
\[
  T^\mu{}_\mu(x)=\sum_I \beta_I\,\mathcal O_I(x).
\]
The coefficients \(\beta_I\) are the beta functions.

In a Lagrangian description, the \(\mathcal O_I(x)\) are constructed as
renormalized versions of bare operators \(\mathcal O_I^{\mathrm{bare}}(x)\)
made out of products of elementary fields and their derivatives.  The
Lagrangian is schematically
\[
  \mathcal L=\sum_I g_I^{\mathrm{bare}}\mathcal O_I^{\mathrm{bare}},
\]
where \(g_I(\mu)\) are renormalized couplings defined in a chosen
renormalization scheme, for example minimal subtraction.  Under
variations of the couplings,
\[
  \delta\mathcal L
  =
  \sum_I \delta g_I^{\mathrm{bare}}\mathcal O_I^{\mathrm{bare}}
  =
  \sum_I \delta g_I(\mu)\,[\mathcal O_I]_\mu .
\]
The running of the renormalized couplings is
\[
  \mu\frac{\dd}{\dd\mu}g_I(\mu)=\beta_I(g(\mu)).
\]
These are the same beta functions which appear in
\[
  T^\mu{}_\mu=\sum_I\beta_I[\mathcal O_I]_\mu,
\]
up to total derivatives.  Thus beta functions control how the
renormalized, physical couplings change with the energy or momentum
scale \(\mu\).  A scale-invariant RG fixed point has
\[
  T^\mu{}_\mu=0.
\]

\subsection{The Ising Fixed Point}

As an example of an RG fixed point, recall the statistical Ising model:
a system of spins on a \(D\)-dimensional lattice \(\Lambda\) in thermal
equilibrium.  For example, in \(D=2\) one may take a square lattice
\[
  x=(i,j),\qquad i,j\in\mathbb Z.
\]
At every site \(x\in\Lambda\) there is a spin \(s_x=\pm1\).  Equivalently
there is a two-dimensional vector space
\[
  V_x=\operatorname{span}\{\ket{\uparrow},\ket{\downarrow}\},
\]
unrelated to the Hilbert space of the underlying QFT.  The lattice spin
system has
\[
  \widetilde{\mathcal H}=\bigotimes_{x\in\Lambda}V_x,
  \qquad
  \widetilde H=-\sum_{\langle x x'\rangle}s_xs_{x'},
\]
where \(\langle x x'\rangle\) denotes neighboring sites.  Observables are
thermal expectation values
\[
  \langle\widehat O\rangle
  =
  \frac1Z\operatorname{Tr}_{\widetilde{\mathcal H}}
  \left(\ee^{-\beta\widetilde H}\widehat O\right),
  \qquad
  Z=\operatorname{Tr}_{\widetilde{\mathcal H}}\ee^{-\beta\widetilde H},
\]
with
\[
  \beta=\frac1T
\]
the inverse temperature.

\begin{center}
\begin{tikzpicture}[scale=0.65,>=Stealth]
  \foreach \i in {0,...,5} {
    \draw[gray!55] (0,\i) -- (5,\i);
    \draw[gray!55] (\i,0) -- (\i,5);
  }
  \foreach \x/\y/\s in {
    0/0/1,1/0/-1,2/0/1,3/0/-1,4/0/1,5/0/1,
    0/1/1,1/1/1,2/1/-1,3/1/1,4/1/-1,5/1/1,
    0/2/-1,1/2/1,2/2/1,3/2/1,4/2/-1,5/2/-1,
    0/3/1,1/3/-1,2/3/1,3/3/-1,4/3/1,5/3/1,
    0/4/-1,1/4/-1,2/4/1,3/4/1,4/4/1,5/4/-1,
    0/5/1,1/5/1,2/5/-1,3/5/1,4/5/-1,5/5/1
  }{
    \ifnum\s=1
      \draw[green!50!black,->,line width=1.1pt] (\x,\y-0.28) -- (\x,\y+0.28);
    \else
      \draw[green!50!black,->,line width=1.1pt] (\x,\y+0.28) -- (\x,\y-0.28);
    \fi
  }
  \node[anchor=west] at (6.0,3.7) {\(x=(i,j)\)};
  \node[anchor=west] at (6.0,3.1) {\(i,j\in\mathbb Z\)};
  \node[anchor=west] at (6.0,1.7) {\(s_x=1\)};
  \draw[green!50!black,->,line width=1.1pt] (6.6,1.0) -- (6.6,1.5);
  \node[anchor=west] at (6.0,0.2) {\(s_x=-1\)};
  \draw[green!50!black,->,line width=1.1pt] (6.6,-0.1) -- (6.6,-0.6);
\end{tikzpicture}
\end{center}

At low temperature there is spontaneous magnetization:
\[
  \langle s\rangle\neq 0.
\]
It turns on at the critical inverse temperature
\[
  \beta_c=\frac1{T_c}.
\]

\begin{center}
\begin{tikzpicture}[scale=0.9,>=Stealth]
  \draw[->] (0,0) -- (5.1,0) node[right] {\(\beta\)};
  \draw[->] (0,0) -- (0,3.0) node[above] {\(\langle s\rangle\)};
  \draw[dashed,gray] (0,2.35) -- (4.9,2.35);
  \draw[very thick,blue] (2.0,0) .. controls (2.05,0.6) and (2.55,1.6) .. (3.4,2.05)
    .. controls (4.0,2.32) and (4.5,2.35) .. (4.85,2.35);
  \node[below,magenta] at (2.0,0) {\(\beta_c=1/T_c\)};
  \node[gray,left] at (0,2.35) {\(1\)};
  \node[above right] at (2.65,2.25) {``spontaneous magnetization''};
\end{tikzpicture}
\end{center}

The spin correlator is
\[
  \langle s_{0,0}s_{N,M}\rangle \equiv f_{N,M}\simeq f(r),
  \qquad
  r=\sqrt{N^2+M^2}.
\]
For \(r\gg 1\) and \(T\) close to \(T_c\), the qualitative behavior is
\[
  f(r)\sim e^{-r/\xi(T)}
  \qquad (T>T_c),
\]
\[
  f(r)\sim r^{-2\Delta}
  \qquad (T=T_c),
\]
while for \(T<T_c\) the correlator tends to a nonzero long-distance
limit set by the magnetization.  The correlation length behaves as
\[
  \xi(T)\sim (T-T_c)^{-\nu}.
\]

\begin{center}
\begin{tikzpicture}[scale=0.8,>=Stealth]
  \draw[->] (0,0) -- (6.3,0) node[right] {\(r\)};
  \draw[->] (0,0) -- (0,3.2) node[above] {\(f(r)\)};
  \draw[thick,green!50!black] (0.4,2.8) .. controls (1.5,2.35) and (3.0,2.0) .. (5.7,1.95);
  \node[green!50!black] at (5.0,2.25) {\(T<T_c\)};
  \draw[thick,blue] (0.4,2.8) .. controls (1.1,2.1) and (2.2,1.25) .. (4.9,0.45);
  \node[blue] at (4.2,0.95) {\(T=T_c,\ f(r)\sim r^{-2\Delta}\)};
  \draw[thick,magenta] (0.4,2.8) .. controls (0.9,1.55) and (1.8,0.55) .. (3.25,0.12);
  \node[magenta] at (2.7,-0.55) {\(f(r)\sim e^{-r/\xi(T)}\)};
\end{tikzpicture}
\end{center}

The critical exponents quoted in the notes are
\[
  D=2:\qquad \Delta=\frac18,\qquad \nu=1,
\]
\[
  D=3:\qquad \Delta=0.518\ldots,\qquad \nu=0.62998\ldots.
\]

The statistical Ising model may be equivalently formulated as a
\(D\)-dimensional Euclidean scalar field theory with a single field
variable \(\phi(x)\), a double-well-like potential \(V(\phi)\),
\(SO(D)\)-breaking higher-derivative kinetic terms, and a
non-Euclidean-invariant UV cutoff.  The key claim is that at
\(T=T_c\), where \(\xi(T_c)=\infty\), the long-distance correlators of
the Ising model are captured by a CFT known as the Ising CFT.
Furthermore,
\[
  \text{Ising CFT}
  =
  \text{IR RG fixed point of \(D\)-dimensional massless \(\phi^4\) theory}
\]
for \(D=2,3\).  Here
\[
  \Delta=\text{scaling dimension of }\phi(x),
\]
and
\[
  D-\frac1\nu=\text{scaling dimension of }\phi^2(x).
\]
These quantities are calculated in the \(\epsilon\)-expansion for
\(D=4-\epsilon\).

\section{Conformal Symmetry}

A CFT is a QFT whose stress-energy tensor \(\widehat T_{\mu\nu}(x)\) is
traceless:
\[
  \widehat T^\mu{}_\mu(x)=0.
\]
The consequence is the existence of more conserved currents, hence more
symmetries.  Consider
\[
  \widehat J_\mu(x)\equiv \widehat T_{\mu\nu}(x)\epsilon^\nu(x)
\]
for some chosen vector field \(\epsilon^\nu(x)\).  Then
\[
  \partial_\mu\widehat J^\mu=0
  \quad\Longleftrightarrow\quad
  \widehat T^{\mu\nu}\partial_\mu\epsilon_\nu=0.
\]
This condition is satisfied provided
\[
  \widehat T^\mu{}_\mu=0
\]
and
\[
  \partial_\mu\epsilon_\nu+\partial_\nu\epsilon_\mu
  -\frac2D\eta_{\mu\nu}\partial_\rho\epsilon^\rho=0.
\]
This is the conformal Killing equation.  The question is: what are the
solutions, namely the conformal Killing vectors?

In \(D>2\), the solutions are
\[
  \epsilon_\mu(x)
  =
  a_\mu+b_{\mu\nu}x^\nu+\lambda x_\mu
  +2(c\cdot x)x_\mu-c_\mu x^2,
\]
where \(a_\mu\), \(c_\mu\), and \(\lambda\) are constants, and
\(b_{\mu\nu}=-b_{\nu\mu}\).  The terms are respectively translations,
Lorentz transformations, dilatations, and special conformal
transformations.  In \(D=2\) there are more conformal Killing vectors;
this case is revisited later.

Given a conformal Killing vector \(\epsilon^\mu(x)\), construct the
conserved Noether current
\[
  \widehat J_\mu=\widehat T_{\mu\nu}\epsilon^\nu.
\]
The corresponding charge is
\[
  \widehat Q_\epsilon
  =
  \int_\Sigma \dd\sigma\,n^\mu\widehat J_\mu,
\]
where \(\Sigma\) is a spacelike slice without boundary and \(n^\mu\) is
its normal timelike vector.  Equivalently, on a time slice,
\[
  \widehat Q_\epsilon
  =
  \int \dd^{D-1}\vec x\,\widehat T_{0\nu}\epsilon^\nu .
\]
It follows from \(\partial_\mu J^\mu=0\) and Stokes' theorem that
\(\widehat Q_\epsilon\) is independent of the choice of \(\Sigma\).

\subsection{Geometric Interpretation}

The vector field \(\epsilon^\mu(x)\) generates an infinitesimal symmetry
transformation
\[
  \widehat U=\ee^{\ii\widehat Q_\epsilon}.
\]
The case
\[
  \epsilon^\mu(x)=a^\mu+b^\mu{}_\nu x^\nu
\]
corresponds to Poincare symmetry.  In particular, such symmetry
transformations compose just like infinitesimal coordinate
transformations
\[
  x^\mu\mapsto x^\mu+\epsilon^\mu(x).
\]
The commutator of charges is expected to extend to all conformal Killing
vectors:
\[
  [\widehat Q_{\epsilon_1},\widehat Q_{\epsilon_2}]
  =
  -\ii\,\widehat Q_{[\epsilon_1,\epsilon_2]},
\]
where the commutator of vector fields is
\[
  [\epsilon_1,\epsilon_2]^\mu
  =
  \epsilon_1^\nu\partial_\nu\epsilon_2^\mu
  -
  \epsilon_2^\nu\partial_\nu\epsilon_1^\mu .
\]
Thus infinitesimal conformal transformations can be labelled by
infinitesimal coordinate transformations, and finite conformal
transformations can be labelled by finite coordinate transformations.

To see what kind of finite transformations arise, write
\[
  \tilde x^\mu=x^\mu+\epsilon^\mu(x).
\]
To first order,
\[
  \frac{\partial\tilde x^\mu}{\partial x^\rho}
  \frac{\partial\tilde x^\nu}{\partial x^\sigma}\eta_{\mu\nu}
  =
  \eta_{\rho\sigma}
  +\partial_\rho\epsilon_\sigma+\partial_\sigma\epsilon_\rho
  +O(\epsilon^2)
  =
  \ee^{2\omega(x)}\eta_{\rho\sigma}
\]
for some \(\omega(x)\).  This relation is preserved under composition.
The conclusion is that finite conformal transformations are coordinate
transformations
\[
  x^\mu\mapsto \tilde x^\mu=f^\mu(x)
\]
such that
\[
  \frac{\partial\tilde x^\mu}{\partial x^\rho}
  \frac{\partial\tilde x^\nu}{\partial x^\sigma}\eta_{\mu\nu}
  =
  \ee^{2\omega(x)}\eta_{\rho\sigma}
\]
for some \(\omega(x)\).

In \(D>2\) dimensions it is easy to write down all finite conformal
transformations.  Consider the inversion
\[
  I:\qquad x^\mu\mapsto \tilde x^\mu=\frac{x^\mu}{x^2},
\]
in Euclidean signature, with
\[
  x^2=\delta_{\mu\nu}x^\mu x^\nu.
\]
This is a conformal transformation, since
\[
  \frac{\partial\tilde x^\mu}{\partial x^\rho}
  =
  \frac1{x^2}
  \left(\delta^\mu{}_\rho-2\frac{x^\mu x_\rho}{x^2}\right),
\]
and therefore it rescales the metric by a position-dependent factor.
The inversion is not continuously connected to the identity.  On the
other hand,
\[
  I\circ \text{translation}\circ I
\]
is continuously connected to the identity and gives a special conformal
transformation,
\[
  x^\mu\mapsto
  \frac{x^\mu+a^\mu x^2}{1+2a\cdot x+a^2x^2}.
\]
Its infinitesimal form is
\[
  x^\mu\mapsto x^\mu+a^\mu x^2-2(a\cdot x)x^\mu+O(a^2),
\]
as seen from the conformal Killing vector.  The Poincare symmetries
together with inversion generate all conformal transformations in
\(D\geq 2\) dimensions.

\subsection{Conformal Versus Weyl Invariance}

A classical field theory action
\[
  S[\phi]=\int \dd^D x\,\mathcal L(\phi,\partial\phi,\ldots)
\]
can be extended to a covariant action \(S[\phi;g]\) for fields
propagating in a curved spacetime metric
\[
  \dd s^2=g_{\mu\nu}(x)\dd x^\mu\dd x^\nu.
\]
For example, for a free scalar in Euclidean signature,
\[
  S_E[\phi]=\int \dd^D x\,
  \left(\frac12\partial_\mu\phi\,\partial^\mu\phi+\frac12m^2\phi^2\right)
\]
has a covariant version of the form
\[
  S_E[\phi;g]
  =
  \int \dd^D x\,\sqrt{\det g}\,
  \left(\frac12 g^{\mu\nu}\partial_\mu\phi\partial_\nu\phi
  +\frac12m^2\phi^2+R(g)\cdots\right).
\]
Under a metric variation,
\[
  S[\phi;g+\delta g]-S[\phi;g]
  =
  \frac12\int \dd^D x\,\sqrt{\det g}\,
  \delta g_{\mu\nu}T^{\mu\nu}.
\]
When restricted to \(g_{\mu\nu}=\eta_{\mu\nu}\), or \(\delta_{\mu\nu}\)
in Euclidean signature, this \(T_{\mu\nu}\) agrees with the Noether
current for translation symmetry.

A local QFT may also be defined in curved spacetime, at least for small
deformations away from flat space,
\[
  g_{\mu\nu}=\eta_{\mu\nu}+\delta g_{\mu\nu},
\]
by inserting
\[
  \exp\left(-\frac12\int \delta g_{\mu\nu}T^{\mu\nu}\right)
\]
in correlation functions.  This should now be interpreted as a local
operator insertion.  There is a caution: there can be potential
singularities when several \(T_{\mu\nu}\) insertions collide beyond first
order in \(\delta g\), due to the OPE as well as contact terms.  This
will be revisited in the context of the Weyl anomaly.

For a CFT,
\[
  T^\mu{}_\mu(x)=0
\]
modulo contact terms and curvature terms.  Thus \(S[\phi;g]\) is
invariant under metric variations
\[
  \delta g_{\mu\nu}\propto g_{\mu\nu}.
\]
Equivalently, \(S[\phi;g]\) is invariant under a Weyl transformation
\[
  g_{\mu\nu}(x)\mapsto \ee^{2\omega(x)}g_{\mu\nu}(x),
\]
which acts on \(g_{\mu\nu}\), and possibly also on \(\phi\).  A CFT is
Weyl invariant when covariantly coupled to a background metric
\(g_{\mu\nu}\).

For the free massless scalar field in \(D\) dimensions, the
covariantized Euclidean action is
\[
  S_0=\int \dd^D x\,\sqrt{\det g}\,
  \frac12 g^{\mu\nu}\partial_\mu\phi\partial_\nu\phi.
\]
Varying with respect to the background metric gives, in flat spacetime,
\[
  T_{\mu\nu}^{(0)}
  =
  \partial_\mu\phi\,\partial_\nu\phi
  -\frac12\eta_{\mu\nu}(\partial\phi)^2,
\]
and hence
\[
  T^\mu{}_\mu
  =
  \left(1-\frac D2\right)(\partial\phi)^2,
\]
which is not zero for \(D>2\).  Thus one asks whether the free massless
scalar theory is a CFT.

Consider the improvement term
\[
  S_1=\int \dd^D x\,\sqrt{\det g}\,R(g)\phi^2 .
\]
Under a metric variation,
\[
  \delta_g S_1
  =
  \frac12\int \dd^D x\,\sqrt{\det g}\,\delta g_{\mu\nu}
  \left(-\partial^\mu\partial^\nu\phi^2
  +\eta^{\mu\nu}\Box\phi^2\right)
\]
in flat spacetime, up to terms which vanish for \(g_{\mu\nu}=\eta_{\mu\nu}\).
Although \(S_1\) itself vanishes in flat spacetime, \(\delta_g S_1\)
contributes to \(T_{\mu\nu}\) in flat spacetime.

Let
\[
  S=S_0+aS_1 .
\]
Then in flat spacetime
\[
  T_{\mu\nu}
  =
  \partial_\mu\phi\,\partial_\nu\phi
  -\frac12\eta_{\mu\nu}(\partial\phi)^2
  +a\left(-\partial_\mu\partial_\nu\phi^2+\eta_{\mu\nu}\Box\phi^2\right).
\]
The improvement term does not contribute to
\[
  \widehat P^\mu=\int \dd^{D-1}\vec x\,T^{0\mu},
\]
but it does contribute to the trace:
\[
\begin{aligned}
  T^\mu{}_\mu
  &=
  \left(1-\frac D2\right)(\partial\phi)^2
  +a(D-1)\Box\phi^2
\\
  &=
  \left(1-\frac D2+2(D-1)a\right)(\partial\phi)^2
  +2(D-1)a\,\phi\,\Box\phi.
\end{aligned}
\]
Using the equation of motion \(\Box\phi=0\), this vanishes if and only if
\[
  a=\frac{D-2}{4(D-1)}.
\]
In this case
\[
  S=
  \int \dd^D x\,\sqrt{\det g}\,
  \left(
    \frac12 g^{\mu\nu}\partial_\mu\phi\partial_\nu\phi
    +\frac{a}{2}R(g)\phi^2
  \right)
\]
is Weyl invariant, in the sense that \(S\) is invariant under
\[
  \delta g_{\mu\nu}=2\delta\omega\,g_{\mu\nu}
\]
and a suitable choice of \(\delta\phi\).  The need to transform
\(\phi\) is due to the term proportional to the equation of motion in
\(T^\mu{}_\mu\).  This is the conformally coupled scalar.  The statement
\[
  T^\mu{}_\mu(x)=0
\]
then holds as an operator equation in the corresponding free massless
scalar QFT.

\section{Basic Consequences of Conformal Symmetry}

First, conformal invariance includes scale invariance, namely
dilatation invariance.  The corresponding Noether current is
\[
  S_\mu=T_{\mu\nu}x^\nu.
\]
This is possible only when there is no characteristic mass scale
intrinsic to the theory.  For instance, four-dimensional Yang--Mills
theory is classically conformally invariant, but quantum Yang--Mills
theory has a characteristic mass scale \(\Lambda_{\mathrm{QCD}}\), as
well as non-vanishing beta functions, hence
\[
  T^\mu{}_\mu\neq0.
\]
It is therefore not a CFT.  In such QFTs there is no obvious notion of
asymptotic states or particles, and no obvious \(S\)-matrix.

Second, the physical observables of a CFT include correlation functions
of local operators
\[
  \widehat{\mathcal O}_I(x)
\]
which obey microcausality:
\[
  [\widehat{\mathcal O}_I(x),\widehat{\mathcal O}_J(y)]=0
  \qquad\text{for}\qquad (x-y)^2>0,
\]
with \(x,y\in\mathbb R^{1,D-1}\).

There are several versions of correlators.  The Wightman functions are
\[
  \bra{\Omega}
  \widehat{\mathcal O}_{I_1}(x_1)\cdots
  \widehat{\mathcal O}_{I_n}(x_n)
  \ket{\Omega}.
\]
The time-ordered Green functions are
\[
  \bra{\Omega}
  T\{\widehat{\mathcal O}_{I_1}(x_1)\cdots
  \widehat{\mathcal O}_{I_n}(x_n)\}
  \ket{\Omega}.
\]
The Euclidean Green functions are
\[
  \left\langle
  \mathcal O_{I_1}(x_1^E)\cdots \mathcal O_{I_n}(x_n^E)
  \right\rangle .
\]
Wightman and Euclidean correlators are related by analytic continuation
\[
  x_i^0=-\ii x_i^D,
  \qquad
  x_i^E=(x_i^1,\ldots,x_i^D),
\]
while maintaining the ordering
\[
  \operatorname{Im}x_1^0<
  \operatorname{Im}x_2^0<
  \cdots<
  \operatorname{Im}x_n^0.
\]

\begin{center}
\begin{tikzpicture}[scale=0.8,>=Stealth]
  \draw[->] (-0.3,0) -- (5.2,0) node[right] {\(\operatorname{Re}x^0\)};
  \draw[->] (0,-0.3) -- (0,3.3) node[above] {\(-\operatorname{Im}x^0\)};
  \foreach \x/\y/\lab in {0.9/0.45/{x_1},1.7/1.0/{x_2},2.6/1.55/{x_3},3.7/2.3/{x_n}} {
    \fill (\x,\y) circle (2pt);
    \node[above right] at (\x,\y) {\(\lab\)};
  }
  \draw[->,blue,thick] (4.35,0.35) arc[start angle=-20,end angle=70,radius=0.9];
  \node[blue] at (4.55,1.55) {Wick rotation};
\end{tikzpicture}
\end{center}

In the path-integral formulation the Wightman function is obtained from
the Lorentzian path integral
\[
  \bra{\Omega}
  \widehat{\mathcal O}_{I_1}(x_1)\cdots
  \widehat{\mathcal O}_{I_n}(x_n)
  \ket{\Omega}
  =
  \int [D\phi]\,\ee^{\ii S[\phi]}\,
  \mathcal O_{I_1}(x_1)\cdots \mathcal O_{I_n}(x_n),
\]
with the insertion positions placed in the appropriate complex-time
order.  The time-ordered correlator and the Euclidean correlator are
obtained by the corresponding contour deformation, namely Wick
rotation.

\section{State/Operator Correspondence}

Assume that we can covariantize a CFT, for example by a Weyl-invariant
action
\[
  S[\phi;g],
\]
where \(g\) is a background metric.  Starting with Euclidean spacetime
\(\mathbb R^D\), write the metric in polar coordinates as
\[
  \dd s^2=\sum_{\mu=1}^D \dd x^\mu \dd x^\mu
  =
  \dd r^2+r^2\dd\Omega^2,
\]
where \(\dd\Omega^2\) is the line element on the unit sphere
\(S^{D-1}\).  There is a Weyl transformation which changes the
background metric to
\[
  \frac1{r^2}\dd s^2
  =
  \frac{\dd r^2}{r^2}+\dd\Omega^2
  =
  \dd\tau^2+\dd\Omega^2,
  \qquad
  \tau\equiv \log r .
\]
This is the metric of
\[
  \mathbb R\times S^{D-1},
\]
the Euclidean cylinder.  Thus
\[
  \mathbb R^D
  \quad \xleftrightarrow{\text{Weyl}} \quad
  \mathbb R\times S^{D-1}.
\]

\begin{center}
\begin{tikzpicture}[scale=0.9,>=Stealth]
  \draw[thick] (-5.2,-2.2) rectangle (-1.1,2.0);
  \node at (-5.0,2.25) {\(\mathbb R^D\)};
  \draw[green!45!black,thick] (-3.15,-0.25) ellipse (1.45 and 1.05);
  \draw[green!45!black,thick,dashed] (-4.45,-0.25) arc[start angle=180,end angle=360,x radius=1.3,y radius=0.34];
  \draw[blue,->,line width=1pt] (-3.15,-0.25) -- (-2.25,0.95) node[above] {\(r\)};
  \draw[blue,->,line width=1pt] (-3.15,-0.25) -- (-3.15,1.25);
  \draw[blue,->,line width=1pt] (-3.15,-0.25) -- (-1.95,-0.75);
  \fill[purple] (-3.15,-0.25) circle (2pt);
  \node[purple] at (-3.4,-0.05) {\(\mathcal O(0)\)};
  \fill[purple] (-2.1,-1.65) circle (2pt);
  \node[purple,right] at (-2.0,-1.55) {\(\mathcal O(x)\)};

  \draw[<->,thick] (-0.3,0) -- (1.2,0);
  \node[below] at (0.45,-0.15) {\(r=\ee^\tau\)};

  \draw[thick] (2.0,-2.3) -- (2.0,2.1);
  \draw[thick] (5.2,-2.3) -- (5.2,2.1);
  \draw[thick] (2.0,2.1) ellipse (1.6 and 0.32);
  \draw[thick] (2.0,-2.3) arc[start angle=180,end angle=360,x radius=1.6,y radius=0.32];
  \draw[thick,dashed] (2.0,-2.3) arc[start angle=180,end angle=0,x radius=1.6,y radius=0.32];
  \draw[green!45!black,thick] (2.0,-0.45) arc[start angle=180,end angle=360,x radius=1.6,y radius=0.32];
  \draw[green!45!black,thick,dashed] (2.0,-0.45) arc[start angle=180,end angle=0,x radius=1.6,y radius=0.32];
  \foreach \x in {2.7,3.65,4.55} {
    \draw[blue,->,line width=1pt] (\x,-1.15) -- (\x,-0.35);
  }
  \draw[->,thick] (1.6,-2.25) -- (1.6,2.0) node[above] {\(\mathbb R\)};
  \node at (1.35,0.7) {\(\tau\)};
  \node at (3.6,-2.8) {\(\Omega\in S^{D-1}\)};
  \fill[purple] (3.85,-0.05) circle (2pt);
  \node[purple,right] at (3.95,0.05) {\(\widetilde{\mathcal O}(\tau,\Omega)\)};
  \node[purple] at (2.0,-2.75) {\(\ket{\mathcal O}\)};
\end{tikzpicture}
\end{center}

What does this equivalence mean?  The CFT on \(\mathbb R^D\), characterized
by the data of its spectrum of local operators and their Euclidean
correlation functions, is equivalent to a CFT on
\(\mathbb R\times S^{D-1}\), with suitable curvature coupling, with the
same spectrum of local operators and identical correlation functions up
to operator redefinition.  In formulas,
\[
  \left\langle
    \mathcal O_1(x_1)\cdots\mathcal O_n(x_n)
  \right\rangle_{\mathbb R^D}
  =
  \left\langle
    \widetilde{\mathcal O}_1(\tau_1,\Omega_1)\cdots
    \widetilde{\mathcal O}_n(\tau_n,\Omega_n)
  \right\rangle_{\mathbb R\times S^{D-1}}.
\]
We will understand the map
\[
  \mathcal O\longmapsto \widetilde{\mathcal O}
\]
later.

The Weyl transformation is singular at \(r=0\), namely
\(\tau=-\infty\).  Hence an operator insertion at the origin,
\[
  \mathcal O(x=0),
\]
is mapped to a state
\[
  \ket{\mathcal O}
\]
of the CFT on \(\mathbb R\times S^{D-1}\), where \(\mathbb R\) is
Euclidean time and \(S^{D-1}\) is space.  The expected relation is
\[
  \left\langle
    \mathcal O(0)\mathcal O_1(x_1)\cdots\mathcal O_n(x_n)
  \right\rangle_{\mathbb R^D}
  =
  \left\langle
    \Omega\middle|
    \widetilde{\mathcal O}_1(\tau_1,\Omega_1)\cdots
    \widetilde{\mathcal O}_n(\tau_n,\Omega_n)
    \middle|\mathcal O
  \right\rangle_{\mathbb R\times S^{D-1}},
\]
assuming
\[
  \tau_1>\tau_2>\cdots>\tau_n .
\]
Here \(\ket{\Omega}\) denotes the ground state on \(S^{D-1}\).

Scaling transformations correspond to the conformal Killing vector
\[
  r\frac{\partial}{\partial r}
  \quad\Longleftrightarrow\quad
  \frac{\partial}{\partial\tau},
\]
namely time translation on \(\mathbb R\times S^{D-1}\).  Therefore
dilatation on \(\mathbb R^D\) corresponds to the Hamiltonian on
\(S^{D-1}\).  We will elaborate on the meaning of this later.

\subsection{Conformally Coupled Free Scalar on the Cylinder}

For the conformally coupled free scalar,
\[
  S=
  \int \dd^D x\,\sqrt{\det g}\,
  \left(
    \frac12 g^{\mu\nu}\partial_\mu\phi\,\partial_\nu\phi
    +\frac12 a R(g)\phi^2
  \right),
  \qquad
  a=\frac{D-2}{4(D-1)} .
\]
On \(\mathbb R\times S^{D-1}\),
\[
  R(g)=(D-1)(D-2).
\]
Thus the theory becomes
\[
  \widetilde S=
  \int \dd^D x\,\sqrt{\det g}\,
  \left(
    \frac12(\partial\phi)^2+\frac12 m^2\phi^2
  \right),
\]
with conformal mass
\[
  m=\sqrt{aR}=\frac{D-2}{2}.
\]
This is the same value as the mass or scaling dimension of the field
operator \(\widehat\phi\) on \(\mathbb R^D\).

An exercise is to start with the Euclidean action of a free massless
scalar field \(S_0[\phi]\) on \(\mathbb R^D\), perform the coordinate
change
\[
  x^\mu\mapsto (r,\Omega),
\]
together with a suitable coordinate-dependent rescaling
\[
  \phi(x)\mapsto \widetilde\phi(r,\Omega),
\]
and show that
\[
  S_0[\phi]=S[\widetilde\phi]
\]
on \(\mathbb R\times S^{D-1}\), with the conformal mass term.

\subsection{Local Operators of the Free Scalar}

What are the local operators?  They are
\[
  \phi,\qquad
  \partial_\mu\phi,\qquad
  \partial_\mu\partial_\nu\phi,\qquad \ldots
\]
and their products, which must be regularized.  For example,
\[
  :\phi(x)^2:
  \equiv
  \lim_{y\to x}
  \left(
    \phi(x)\phi(y)-\left\langle\phi(x)\phi(y)\right\rangle
  \right),
\]
where the subtracted two-point function is the Green function
\[
  G(x,y)=\left\langle\phi(x)\phi(y)\right\rangle .
\]
This makes correlation functions such as
\[
  \left\langle
    :\phi(x)^2:\,\phi(y)\phi(z)\cdots
  \right\rangle
\]
well-defined.

There is also a redundancy: by the equation of motion,
\[
  \Box\phi(x)=0.
\]
Thus \(\Box\phi(x)\) vanishes as a local operator, and so do
\[
  :\phi\Box\phi:,\qquad
  :(\Box\phi)^2:,\qquad \ldots
\]

It is illuminating to enumerate all linearly independent local
operators using an operator partition function
\[
  \mathcal V_{\mathrm{loc}}
  =
  \text{space of all local operators at a fixed point},
\]
\[
  Z(\beta)
  \equiv
  \operatorname{Tr}_{\mathcal V_{\mathrm{loc}}}q^\Delta,
  \qquad
  q=\ee^{-\beta},
\]
where \(\Delta\) is the scaling dimension.  For example,
\[
  \Delta_\phi=\frac{D-2}{2},
  \qquad
  \Delta_{\partial_\mu\phi}=\frac{D-2}{2}+1,
  \qquad
  \Delta_{\partial_\mu\partial_\nu\phi}=\frac{D-2}{2}+2,
  \qquad \ldots
\]

A complete basis of local operators in the free massless scalar field
theory consists of unordered words
\[
  :\phi\,\partial_\mu\phi\,\partial_\nu\phi\cdots:
\]
made from letters, subject to \(\Box\phi=0\).  The letter partition
function, counting
\(\phi,\partial_\mu\phi,\ldots\), is
\[
\begin{aligned}
  f(q)
  &\equiv
  \operatorname{Tr}_{\mathrm{letters}} q^\Delta
\\
  &=
  q^{\frac{D-2}{2}}\,
  \frac{1-q^2}{(1-q)^D}
\\
  &=
  \sum_a q^{\Delta_a}.
\end{aligned}
\]
The factor \(q^{(D-2)/2}\) is the contribution of \(\Delta_\phi\), the
denominator counts the \(D\) different derivatives
\(\partial_\mu\), and the factor \(1-q^2\) removes \(\Box\phi\).

The partition function of unordered words is related by
\[
  Z(\beta)
  =
  \prod_a \frac1{1-q^{\Delta_a}}
  =
  \exp\left(
    \sum_{n=1}^{\infty}\frac1n f(q^n)
  \right).
\]

The state/operator map would identify each local operator with a state
of the conformally coupled scalar field theory on \(S^{D-1}\).  For
example, the single-particle partition function is
\[
  \operatorname{Tr}_{\text{1-particle}}\ee^{-\beta \widehat H},
\]
where this is a relativistic particle of mass \(m\) with Hamiltonian
\[
  \widehat H=\sqrt{-\nabla^2+m^2}
\]
acting on wavefunctions
\[
  \psi\in L^2(S^{D-1}).
\]
For \(D=3\), the space \(L^2(S^2)\) is spanned by spherical harmonics
labelled by
\[
  \ell=0,1,2,\ldots,\qquad m=-\ell,-\ell+1,\ldots,\ell,
\]
with
\[
  -\nabla^2=\ell(\ell+1).
\]
The conformal mass is
\[
  m=\frac{D-2}{2}=\frac12,
\]
so
\[
  \widehat H
  =
  \sqrt{\ell(\ell+1)+\frac14}
  =
  \ell+\frac12.
\]
Therefore the single-particle partition function is
\[
  \sum_{\ell=0}^{\infty}(2\ell+1)\ee^{-\beta(\ell+1/2)}
  =
  f(\ee^{-\beta})
  =
  \frac{\ee^{-\beta/2}(1+\ee^{-\beta})}
  {(1-\ee^{-\beta})^2},
\]
matching the letter partition function for \(D=3\).

\section{Ward Identities and Conformal Charges}

In quantum mechanics, given a symmetry generator \(\widehat Q\) and an
operator \(\widehat{\mathcal O}\), the infinitesimal symmetry variation
of the operator is
\[
  \delta_\epsilon \widehat{\mathcal O}
  =
  \ii\epsilon[\widehat Q,\widehat{\mathcal O}] .
\]
In quantum field theory, \(\widehat Q\) is obtained by integrating the
Noether current along a spatial slice:
\[
  \widehat Q(\Sigma)
  =
  \int_\Sigma *\widehat j
  =
  \int_\Sigma \dd\sigma\, n^\mu \widehat j_\mu .
\]
The same statement also applies to imaginary time, with the usual Wick
rotation factor.

\begin{center}
\begin{tikzpicture}[scale=0.86,>=Stealth]
  \draw[->] (-5.6,-1.9) -- (-5.6,2.0) node[above] {time};
  \draw[->] (-5.8,-1.65) -- (-3.0,-1.65) node[right] {space};

  \draw[green!45!black,thick]
    (-5.0,-1.0) -- (-2.1,-0.75) -- (-1.2,-0.25) -- (-4.1,-0.45) -- cycle;
  \draw[green!45!black,thick]
    (-4.7,0.75) -- (-1.8,0.95) -- (-0.9,1.45) -- (-3.8,1.25) -- cycle;
  \node[green!45!black] at (-1.3,1.85) {\(\Sigma_2\)};
  \node[green!45!black] at (-1.55,-0.8) {\(\Sigma_1\)};
  \fill[red] (-3.2,0.1) circle (2pt);
  \node[red,below right] at (-3.15,0.1) {\(\widehat{\mathcal O}(0)\)};

  \draw[->,blue,thick] (-0.35,0.35) -- (1.15,0.35)
    node[midway,above] {deform};

  \draw[green!45!black,thick]
    (2.2,-0.15) .. controls (2.25,1.2) and (3.2,1.7) .. (4.5,1.45)
    .. controls (5.25,1.25) and (5.45,0.5) .. (5.05,-0.25)
    .. controls (4.55,-1.2) and (3.0,-1.1) .. (2.45,-0.6)
    .. controls (2.25,-0.45) and (2.15,-0.3) .. (2.2,-0.15);
  \draw[green!45!black,thick,dashed]
    (2.55,-0.45) .. controls (3.2,-0.75) and (4.15,-0.72) .. (4.85,-0.3);
  \node[green!45!black] at (5.55,0.35) {\(\Sigma\)};
  \fill[red] (3.65,0.15) circle (2pt);
  \node[red,above right] at (3.68,0.15) {\(\widehat{\mathcal O}(0)\)};
\end{tikzpicture}
\end{center}

Thus
\[
  [\widehat Q,\widehat{\mathcal O}(0)]
  =
  \widehat Q(\Sigma_2)\widehat{\mathcal O}(0)
  -
  \widehat{\mathcal O}(0)\widehat Q(\Sigma_1)
  =
  T\,\widehat Q(\Sigma)\widehat{\mathcal O}(0),
\]
where \(T\) denotes time-ordering.  Viewed as local operators inserted
in a Euclidean Green function, we may write
\[
  [Q,\mathcal O(0)]=Q(\Sigma)\mathcal O(0),
\]
where \(\Sigma\) is a closed hypersurface in \(\mathbb R^D\) that
encloses the origin.

The implication for correlation functions is the Ward identity.  If the
vacuum is invariant under the symmetry,
\[
  \widehat Q\ket{\Omega}=0,
\]
then
\[
\begin{aligned}
  &\bra{\Omega}
  \delta_\epsilon\widehat{\mathcal O}_1(x_1)
  \widehat{\mathcal O}_2(x_2)\cdots
  \widehat{\mathcal O}_n(x_n)
  \ket{\Omega}
\\
  &\quad+
  \bra{\Omega}
  \widehat{\mathcal O}_1(x_1)
  \delta_\epsilon\widehat{\mathcal O}_2(x_2)\cdots
  \widehat{\mathcal O}_n(x_n)
  \ket{\Omega}
  +\cdots
\\
  &\quad+
  \bra{\Omega}
  \widehat{\mathcal O}_1(x_1)\cdots
  \widehat{\mathcal O}_{n-1}(x_{n-1})
  \delta_\epsilon\widehat{\mathcal O}_n(x_n)
  \ket{\Omega}
  =0 .
\end{aligned}
\]
Equivalently, for Euclidean Green functions,
\[
  \sum_{i=1}^n
  \left\langle
    \mathcal O_1(x_1)\cdots
    \delta_\epsilon\mathcal O_i(x_i)\cdots
    \mathcal O_n(x_n)
  \right\rangle
  =0.
\]
In terms of small hypersurfaces \(\Sigma_i\) enclosing \(x_i\), this is
the statement
\[
  \sum_{i=1}^n
  \left\langle
    \mathcal O_1(x_1)\cdots
    \ii\epsilon\,Q(\Sigma_i)\mathcal O_i(x_i)\cdots
    \mathcal O_n(x_n)
  \right\rangle
  =0.
\]

\begin{center}
\begin{tikzpicture}[scale=0.72,>=Stealth]
  \draw[thick] (-4.8,-2.6) rectangle (4.8,2.7);
  \node at (-4.55,2.95) {\(\mathbb R^D\)};

  \draw[green!45!black,thick]
    (-3.8,-1.3) .. controls (-4.3,0.3) and (-3.0,2.2) .. (-1.1,2.15)
    .. controls (0.5,2.5) and (3.1,1.9) .. (3.55,0.35)
    .. controls (4.1,-1.55) and (1.7,-2.15) .. (-0.5,-1.95)
    .. controls (-2.0,-1.8) and (-3.45,-2.05) .. (-3.8,-1.3);
  \node[green!45!black] at (-0.2,-2.35) {\(\Sigma\)};

  \foreach \x/\y/\lab/\surf in {
    -2.1/0.75/{\mathcal O(x_1)}/{\Sigma_1},
     0.75/0.95/{\mathcal O(x_2)}/{\Sigma_2},
     2.15/-0.75/{\mathcal O(x_n)}/{\Sigma_n}
  }{
    \draw[blue,thick] (\x,\y) ellipse (0.78 and 0.48);
    \draw[blue,thick,dashed] (\x-0.72,\y) arc[start angle=180,end angle=360,x radius=0.72,y radius=0.16];
    \fill (\x-0.18,\y-0.1) circle (1.8pt);
    \node at (\x+0.05,\y+0.18) {\(\lab\)};
    \node[blue] at (\x+0.65,\y+0.56) {\(\surf\)};
  }
  \node at (1.35,-0.15) {\(\cdots\)};
  \node[right] at (5.25,0.85) {\(\Sigma_1\cup\Sigma_2\cup\cdots\cup\Sigma_n\)};
  \draw[->,thick] (5.65,0.45) -- (5.65,-0.95);
  \node[right] at (5.25,-1.25) {deform to \(\Sigma\)};
\end{tikzpicture}
\end{center}

The symmetry invariance of the correlator amounts to
\[
  \left\langle
    Q(\Sigma)\mathcal O_1(x_1)\cdots\mathcal O_n(x_n)
  \right\rangle
  =0,
\]
where \(\Sigma\) encloses all of \(x_1,\ldots,x_n\).  These are the Ward
identities.

Given a conformal Killing vector \(\epsilon^\mu(x)\), the corresponding
conformal symmetry generator is
\[
  Q_\epsilon
  =
  \int_\Sigma \dd\sigma\, n^\mu T_{\mu\nu}(x)\epsilon^\nu(x).
\]

\subsection{Conformal Generators on the Cylinder}

It is illuminating to inspect the conformal Killing vectors under the
operator/state map
\[
  |x|=\ee^\tau,
  \qquad
  \mathbb R^D\longleftrightarrow \mathbb R\times S^{D-1},
  \qquad
  \mathcal O(0)\longleftrightarrow
  \ket{\mathcal O}\in\mathcal H_{S^{D-1}} .
\]
The inversion
\[
  I:\quad x^\mu\mapsto \frac{x^\mu}{x^2}
\]
corresponds to time reversal on the cylinder:
\[
  (\tau,\Omega)\mapsto (-\tau,\Omega).
\]
The dilatation generator
\[
  \widehat D:\quad x^\mu\partial_\mu
\]
corresponds to time translation,
\[
  \partial_\tau,
\]
and the rotation generator
\[
  \widehat J_{\mu\nu}:\quad
  x_\mu\partial_\nu-x_\nu\partial_\mu
\]
corresponds to the \(SO(D)\) rotations of \(S^{D-1}\).

Writing \(x^\mu=r n^\mu\), with \(n^\mu n_\mu=1\), the translation
generator is
\[
\begin{aligned}
  \widehat P_\mu:\quad
  \partial_\mu
  &=
  n_\mu\frac{\partial}{\partial r}
  +
  \frac{\delta_{\mu\nu}-n_\mu n_\nu}{r}
  \frac{\partial}{\partial n_\nu}
\\
  &=
  \ee^{-\tau}
  \left(
    n_\mu\partial_\tau
    +
    (\delta_{\mu\nu}-n_\mu n_\nu)
    \frac{\partial}{\partial n_\nu}
  \right).
\end{aligned}
\]
The special conformal generator is
\[
\begin{aligned}
  \widehat K_\mu:\quad
  I\partial_\mu I
  &=
  x^2\partial_\mu-2x_\mu x^\nu\partial_\nu
\\
  &=
  \ee^\tau
  \left(
    -n_\mu\partial_\tau
    +
    (\delta_{\mu\nu}-n_\mu n_\nu)
    \frac{\partial}{\partial n_\nu}
  \right).
\end{aligned}
\]
On the Lorentzian cylinder, parameterized by \((t,\Omega)\) with
\[
  t=-\ii\tau,
\]
the generators \(\widehat P_\mu\) and \(\widehat K_\mu\) are associated
with complex conformal Killing vectors that are complex conjugates of
one another.  Hence, as operators on \(\mathcal H_{S^{D-1}}\),
\[
  \widehat P_\mu^\dagger=\widehat K_\mu .
\]

Let us inspect the dilatation generator more explicitly, as a linear
operator acting on local operators \(\mathcal O(0)\) at the origin of
\(\mathbb R^D\).  It can be constructed as
\[
\begin{aligned}
  \widehat D
  &=
  -\ii
  \int_{S_r^{D-1}}\dd\sigma\,
  n^\mu \widehat T_{\mu\nu}(x)x^\nu
\\
  &=
  -\ii r^D
  \int_{S^{D-1}}\dd\Omega\,
  n^\mu n^\nu\widehat T_{\mu\nu}(r\vec n),
\end{aligned}
\]
which is independent of \(r\) by the conservation law.

\begin{center}
\begin{tikzpicture}[scale=0.8,>=Stealth]
  \draw[thick] (0,0) ellipse (2.25 and 1.15);
  \draw[thick,dashed] (-2.25,0) arc[start angle=180,end angle=360,x radius=2.25,y radius=0.35];
  \draw[thick] (-2.25,0) arc[start angle=180,end angle=0,x radius=2.25,y radius=0.35];
  \fill (0,0) circle (2pt);
  \node[below left] at (0,0) {\(0\)};
  \coordinate (p) at (1.15,0.72);
  \fill[blue] (p) circle (2pt);
  \draw[blue,->,line width=1pt] (0,0) -- (p) node[midway,above left] {\(r\)};
  \draw[green!45!black,->,line width=1pt] (p) -- ++(0.85,0.55)
    node[above right] {\(\vec n\)};
  \node at (2.9,0.15) {\(S_r^{D-1}\)};
\end{tikzpicture}
\end{center}

The Hamiltonian of the CFT on \(S^{D-1}\) space, on the other hand, is
\[
  \widehat H
  =
  \int_{S^{D-1}}\dd\Omega\,\widehat T_{tt}
  =
  -\int_{S^{D-1}}\dd\Omega\,\widehat T_{\tau\tau},
\]
where \(\tau\) is the Euclidean time coordinate on
\(\mathbb R\times S^{D-1}\).  The identification is
\[
  \widehat D\big|_{\mathbb R^D}
  \longleftrightarrow
  \ii\widehat H\big|_{\text{cylinder}},
\]
and
\[
  r^D\widehat T_{\mu\nu}(r\vec n)n^\mu n^\nu\big|_{\mathbb R^D}
  \longleftrightarrow
  \widehat T_{\tau\tau}(\tau,\Omega)\big|_{\text{cylinder}}.
\]
The identification between the stress-energy tensor on \(\mathbb R^D\)
and on \(\mathbb R\times S^{D-1}\) involves more than just the coordinate
transformation of a tensor; it also involves a suitable weight factor.

With the identification
\[
  \text{``space of local operators at the origin''}
  \equiv
  \mathcal H_{S^{D-1}},
\]
we will think of the conformal charges \(Q_\epsilon\) either as linear
operations on local operators \(\mathcal O(0)\), or as operators on
\(\mathcal H_{S^{D-1}}\).

\subsection{The Conformal Algebra}

The basic conformal charges and the associated conformal Killing
vectors are
\[
\begin{array}{c|c}
  \text{charge} & \text{conformal Killing vector} \\
  \hline
  \widehat P_\mu & \partial_\mu \\
  \widehat J_{\mu\nu} & x_\mu\partial_\nu-x_\nu\partial_\mu \\
  \widehat D & x^\mu\partial_\mu \\
  \widehat K_\mu & x^2\partial_\mu-2x_\mu x^\nu\partial_\nu
\end{array}
\]
Their commutation relations are encoded by
\[
  [\widehat Q_{\epsilon_1},\widehat Q_{\epsilon_2}]
  =
  -\ii \widehat Q_{[\epsilon_1,\epsilon_2]} .
\]
Explicitly,
\[
  [\widehat P_\mu,\widehat J_{\nu\rho}]
  =
  -\ii\left(
    \eta_{\mu\nu}\widehat P_\rho
    -(\nu\leftrightarrow\rho)
  \right),
\]
\[
  [\widehat J_{\mu\nu},\widehat J_{\rho\sigma}]
  =
  -\ii\left(
    \eta_{\nu\rho}\widehat J_{\mu\sigma}
    -
    \eta_{\mu\rho}\widehat J_{\nu\sigma}
    -(\rho\leftrightarrow\sigma)
  \right),
\]
and
\[
  [\widehat D,\widehat P_\mu]=\ii\widehat P_\mu,\qquad
  [\widehat D,\widehat K_\mu]=-\ii\widehat K_\mu,\qquad
  [\widehat D,\widehat J_{\mu\nu}]=0,
\]
\[
  [\widehat P_\mu,\widehat K_\nu]
  =
  2\ii\left(
    \eta_{\mu\nu}\widehat D-\widehat J_{\mu\nu}
  \right).
\]
Together they generate the conformal algebra, as a Lie algebra
isomorphic to \(\mathfrak{so}(D,2)\).  Define
\[
  \widehat{\mathcal J}_{\mu\nu}
  \equiv
  \widehat J_{\mu\nu},
  \qquad
  \widehat{\mathcal J}_{D,-1}
  \equiv
  \widehat D,
\]
\[
  \widehat{\mathcal J}_{D,\mu}
  \equiv
  \frac12(\widehat P_\mu-\widehat K_\mu),
  \qquad
  \widehat{\mathcal J}_{-1,\mu}
  \equiv
  \frac12(\widehat P_\mu+\widehat K_\mu).
\]
Then
\[
  [\widehat{\mathcal J}_{MN},\widehat{\mathcal J}_{PQ}]
  =
  -\ii\left(
    \eta_{NP}\widehat{\mathcal J}_{MQ}
    -
    \eta_{MP}\widehat{\mathcal J}_{NQ}
    -(P\leftrightarrow Q)
  \right),
\]
where
\[
  M,N,P,Q=-1,0,1,\ldots,D-1,D,
\]
and
\[
  \eta_{MN}=\operatorname{diag}(-1,-1,1,\ldots,1).
\]

States of the CFT on \(S^{D-1}\) are organized in representations of the
conformal algebra:
\[
  \mathcal H_{S^{D-1}}
  \simeq
  \bigoplus_i V_i^{\mathrm{irr}},
\]
where \(V_i^{\mathrm{irr}}\) is an irreducible representation of the
conformal algebra.  Each irreducible representation is generated by a
lowest-energy state \(\ket{\phi_i}\), which by the state/operator map
also corresponds to a local operator \(\widehat\phi_i(0)\).  Such a state
or operator \(\phi_i\) is called a conformal primary state or operator.

\subsection{Conformal Primary States}

Suppose \(\ket{\phi}\) is a conformal primary, which is also assumed to
be an energy eigenstate,
\[
  \widehat H\ket{\phi}=\Delta\ket{\phi}.
\]
Since
\[
  \widehat D=\ii\widehat H,
\]
we have
\[
  \widehat D\ket{\phi}=\ii\Delta\ket{\phi}.
\]
Because
\[
  [\widehat D,\widehat K_\mu]=-\ii\widehat K_\mu,
\]
the operator \(\widehat K_\mu\) lowers the energy by \(1\).  Thus
\[
  \widehat K_\mu\ket{\phi}\in V^{\mathrm{irr}},
\]
but \(\ket{\phi}\) is the lowest-energy state in \(V^{\mathrm{irr}}\) by
assumption.  Therefore
\[
  \widehat K_\mu\ket{\phi}=0,
\]
which is equivalent to the definition of a primary.

The operator \(\widehat P_\mu\), on the other hand, raises the energy by
\(1\).  Under the state/operator map,
\[
  \ket{\phi}\longleftrightarrow \widehat\phi(0),
\]
and
\[
  \widehat P_\mu\ket{\phi}
  \longleftrightarrow
  \widehat P_\mu\widehat\phi(0)
  =
  [\widehat P_\mu,\widehat\phi(0)].
\]
Recall that
\[
  \widehat\phi(x+a)
  =
  \ee^{-\ii\widehat P\cdot a}
  \widehat\phi(x)
  \ee^{\ii\widehat P\cdot a},
\]
so
\[
  [\widehat P_\mu,\widehat\phi(x)]
  =
  \ii\partial_\mu\widehat\phi(x).
\]
Therefore
\[
  \widehat P_\mu\ket{\phi}
  =
  \ii\ket{\partial_\mu\phi},
\]
where \(\ket{\partial_\mu\phi}\) is the state that maps to the local
operator \(\partial_\mu\phi(0)\).

Similarly, suppose \(\widehat\phi_\alpha(x)\) obeys the Lorentz
transformation law
\[
  U(\Lambda)\widehat\phi_\alpha(x)U(\Lambda)^{-1}
  =
  (R(\Lambda))_\alpha{}^\beta
  \widehat\phi_\beta(\Lambda x),
\]
where
\[
  U(\Lambda)
  =
  \exp\left(\frac{\ii}{2}\omega_{\mu\nu}\widehat J^{\mu\nu}\right),
  \qquad
  \Lambda^\mu{}_\nu=\delta^\mu{}_\nu+\omega^\mu{}_\nu,
\]
for infinitesimal \(\omega_{\mu\nu}\), and
\[
  R(\Lambda)
  =
  \exp\left(-\frac{\ii}{2}\omega_{\mu\nu}S^{\mu\nu}\right).
\]
The matrices \(S_{\mu\nu}\) obey the same commutation relations as
\(\widehat J_{\mu\nu}\):
\[
  [S_{\mu\nu},S_{\rho\sigma}]
  =
  -\left(
    \eta_{\nu\rho}S_{\mu\sigma}
    -
    \eta_{\mu\rho}S_{\nu\sigma}
    -(\rho\leftrightarrow\sigma)
  \right).
\]
Then
\[
  [\widehat J_{\mu\nu},\widehat\phi_\alpha(0)]
  =
  -(S_{\mu\nu})_\alpha{}^\beta \widehat\phi_\beta(0),
\]
and therefore
\[
  \widehat J_{\mu\nu}\ket{\phi_\alpha}
  =
  -(S_{\mu\nu})_\alpha{}^\beta \ket{\phi_\beta}.
\]

Thus we have determined how a conformal primary transforms under the
conformal algebra:
\[
  \widehat D\ket{\phi}
  =
  \ii\Delta_\phi\ket{\phi},
  \qquad
  \widehat K_\mu\ket{\phi}
  =
  0,
\]
\[
  \widehat P_\mu\ket{\phi}
  =
  \ii\ket{\partial_\mu\phi},
  \qquad
  \widehat J_{\mu\nu}\ket{\phi_\alpha}
  =
  -(S_{\mu\nu})_\alpha{}^\beta\ket{\phi_\beta}.
\]
For a scalar \(\phi\), \(S_{\mu\nu}=0\).  All other states related to
\(\phi\), or to \(\phi_\alpha\) for non-scalar primaries, by the
conformal algebra are descendants obtained by repeatedly acting with
\(\widehat P_\mu\)'s.  For example,
\[
  \widehat P_\mu\ket{\phi},
  \qquad
  \widehat P_\mu\widehat P_\nu\ket{\phi},
  \qquad \ldots
\]
Under the state/operator map,
\[
  \widehat P_\mu\ket{\phi}
  \longleftrightarrow
  \ii\partial_\mu\phi(0),
  \qquad
  \widehat P_\mu\widehat P_\nu\ket{\phi}
  \longleftrightarrow
  \ii\partial_\mu\,\ii\partial_\nu\phi(0),
  \qquad \ldots
\]

\subsection{Local Operators Away from the Origin}

On \(\mathbb R^D\), a local operator at a nonzero point can be viewed as
a limit of local operators at the origin:
\[
  \widehat{\mathcal O}(x)
  =
  \ee^{x^\mu\partial_\mu}\widehat{\mathcal O}(0)
  =
  \widehat{\mathcal O}(0)
  +x^\mu\partial_\mu\widehat{\mathcal O}(0)
  +\frac12x^\mu x^\nu\partial_\mu\partial_\nu
    \widehat{\mathcal O}(0)+\cdots .
\]
Under the state/operator map,
\[
  \ket{\mathcal O(x)}
  =
  \ee^{-\ii x^\mu \widehat P_\mu}\ket{\mathcal O}.
\]

\begin{figure}[htbp]
\centering
\begin{tikzpicture}[scale=0.9,>=Stealth]
  \draw[thick] (-3.6,-1.6) -- (-3.4,1.65) -- (-0.7,1.75) -- (-0.55,-1.45) -- cycle;
  \node at (-0.25,1.65) {$\mathbb R^D$};
  \draw[orange!85!black,thick]
    (-2.9,-1.05) .. controls (-3.35,0.05) and (-2.85,1.35) .. (-1.95,1.05)
    .. controls (-1.25,0.78) and (-0.95,-0.05) .. (-1.25,-0.85)
    .. controls (-1.75,-1.55) and (-2.45,-1.45) .. cycle;
  \draw[orange!85!black,thick]
    (-2.9,-0.95) .. controls (-2.0,-0.45) and (-1.55,0.08) .. (-1.25,0.45);
  \draw[blue!70!black,very thick,->] (-2.1,0.2) -- (-1.42,0.45);
  \node[blue!70!black] at (-2.25,0.35) {$\mathcal O(0)$};
  \fill[blue!70!black] (-2.25,-0.1) circle (1.3pt);
  \fill[violet] (-1.22,-0.15) circle (1.3pt);
  \node[violet] at (-0.98,0.02) {$\mathcal O(x)$};
  \node[orange!85!black] at (-1.18,1.22) {$\widehat Q_\epsilon$};
  \draw[<->,thick] (0.45,0.05) -- (1.45,0.05);
  \draw[thick] (2.4,1.45) ellipse (0.75 and 0.18);
  \draw[thick] (1.65,1.45) -- (1.45,-1.25);
  \draw[thick] (3.15,1.45) -- (2.95,-1.25);
  \draw[thick] (2.2,-1.25) ellipse (0.75 and 0.18);
  \draw[orange!85!black,thick]
    (1.50,-0.15) .. controls (1.95,0.20) and (2.58,0.05) .. (3.04,0.32);
  \draw[orange!85!black,dashed,thick]
    (1.50,-0.15) .. controls (1.95,-0.42) and (2.62,-0.50) .. (3.04,0.32);
  \node[orange!85!black] at (3.36,-0.05) {$\widehat Q_\epsilon$};
  \node[right] at (3.1,1.35) {$\mathbb R\times S^{D-1}$};
\end{tikzpicture}
\caption{The same charge surface viewed in flat space and on the cylinder.}
\end{figure}

For an arbitrary operator \(\mathcal O\), not necessarily a primary, an
infinitesimal conformal charge acts as
\[
  \widehat Q_\epsilon\cdot\mathcal O(x)
  \longleftrightarrow
  \widehat Q_\epsilon\ket{\mathcal O(x)}
  =
  -\ii\int_\Sigma\dd\sigma\,n^\mu
  T_{\mu\nu}(y)\epsilon^\nu(y)\cdot \mathcal O(x).
\]
For a primary \(\phi(x)\), the action can instead be deduced from
\(\ket{\phi(x)}=\ee^{-\ii x\cdot\widehat P}\ket{\phi(0)}\) and the
conformal algebra.  For instance,
\[
  \widehat K_\mu\ket{\phi(x)}
  =
  \ee^{-\ii x\cdot\widehat P}
  \ee^{\ii x\cdot\widehat P}
  \widehat K_\mu
  \ee^{-\ii x\cdot\widehat P}\ket{\phi(0)}.
\]
Repeatedly using the commutation relations gives
\[
  \ee^{\ii x\cdot\widehat P}
  \widehat K_\mu
  \ee^{-\ii x\cdot\widehat P}
  =
  \widehat K_\mu
  -2x_\mu x\cdot\widehat P
  +x^2\widehat P_\mu
  +2x^\nu\widehat J_{\mu\nu}
  -2x_\mu\widehat D .
\]
Since \(\widehat K_\mu\ket{\phi(0)}=0\), this becomes, with Lorentz
indices suppressed,
\[
  \widehat K_\mu\ket{\phi(x)}
  =
  \ee^{-\ii x\cdot\widehat P}
  \left(
    -2x_\mu x\cdot\widehat P
    +x^2\widehat P_\mu
    -2x^\nu S_{\mu\nu}
    -2\ii x_\mu\Delta
  \right)\ket{\phi(0)}.
\]

Thus the complete set of infinitesimal conformal transformations of a
primary \(\phi(x)\) is
\[
  \widehat K_\mu\cdot\phi(x)
  =
  \left(
    -2\ii x_\mu x\cdot\partial
    +\ii x^2\partial_\mu
    -2x^\nu S_{\mu\nu}
    -2\ii x_\mu\Delta
  \right)\phi(x),
\]
\[
  \widehat P_\mu\cdot\phi(x)
  =
  \ii\partial_\mu\phi(x),
  \qquad
  \widehat D\cdot\phi(x)
  =
  \left(\ii x^\mu\partial_\mu+\ii\Delta\right)\phi(x),
\]
\[
  \widehat J_{\mu\nu}\cdot\phi(x)
  =
  \left(
    \ii(x_\mu\partial_\nu-x_\nu\partial_\mu)-S_{\mu\nu}
  \right)\phi(x).
\]

\subsection{Finite Conformal Transformations of Primaries}

Finite conformal transformations are composed out of infinitesimal ones.
They are labeled by a coordinate map
\[
  x^\mu\longmapsto \widetilde x^\mu(x)
\]
such that
\[
  \delta_{\mu\nu}
  \frac{\partial\widetilde x^\mu}{\partial x^\alpha}
  \frac{\partial\widetilde x^\nu}{\partial x^\beta}
  \propto
  \delta_{\alpha\beta}.
\]
Writing \(U\cdot\phi(x)=\widetilde\phi(x)\), first consider the
infinitesimal case
\[
  \widetilde x^\mu=x^\mu+\epsilon^\mu(x),
  \qquad
  U=\ee^{\ii\widehat Q_\epsilon}.
\]
Then
\[
  U\cdot\phi(x)=\phi(x)+\delta\phi(x),
  \qquad
  \delta\phi(x)=\ii\widehat Q_\epsilon\cdot\phi(x),
\]
where \(\widehat Q_\epsilon\cdot\phi\) is a linear combination of
\(\widehat P\cdot\phi\), \(\widehat K\cdot\phi\), \(\widehat D\cdot\phi\),
and \(\widehat J\cdot\phi\).

For
\[
  \epsilon_\mu(x)
  =
  a_\mu+b_{\mu\nu}x^\nu+\lambda x_\mu
  -2(c\cdot x)x_\mu+c_\mu x^2,
\]
the charge is
\[
  \widehat Q_\epsilon
  =
  a\cdot\widehat P
  -\frac12 b_{\mu\nu}\widehat J^{\mu\nu}
  +\lambda\widehat D
  +c\cdot\widehat K,
\]
and therefore
\[
\begin{aligned}
  \delta\phi(x)
  =
  \Big(
    &-a\cdot\partial
    +b^{\mu\nu}x_\mu\partial_\nu
    +\frac{\ii}{2}b^{\mu\nu}S_{\mu\nu}
    -\lambda x\cdot\partial
    -\lambda\Delta  \\
    &-x^2c\cdot\partial
    +2(c\cdot x)x\cdot\partial
    -2\ii c^\mu x^\nu S_{\mu\nu}
    +2(c\cdot x)\Delta
  \Big)\phi(x).
\end{aligned}
\]
Equivalently,
\[
  \delta\phi(x)
  =
  \left(
    -\epsilon^\mu(x)\partial_\mu
    -\frac{\Delta}{D}\partial_\mu\epsilon^\mu(x)
    -\frac{\ii}{2}\partial_\mu\epsilon_\nu(x)S^{\mu\nu}
  \right)\phi(x).
\]

The finite form of the primary transformation law is
\[
  \widetilde\phi_\alpha(\widetilde x)
  =
  \left(
    \det\frac{\partial\widetilde x^\mu}{\partial x^\nu}
  \right)^{-\Delta/D}
  \bar R_\alpha{}^\beta
  \left(
    \frac{\partial\widetilde x^\mu}{\partial x^\nu}
  \right)
  \phi_\beta(x).
\]
Here
\[
  M^\mu{}_\nu
  \equiv
  \frac{\partial\widetilde x^\mu}{\partial x^\nu}
\]
obeys
\[
  M^T M\propto I
  \qquad\text{(Euclidean)}
\]
or
\[
  M^T\eta M\propto \eta
  \qquad\text{(Lorentzian)}
\]
by the assumption of conformality.  Defining
\[
  \Lambda\equiv (\det M)^{-1/D}M,
\]
we have
\[
  \Lambda^T\Lambda=I
  \qquad\text{(Euclidean)}
\]
or
\[
  \Lambda^T\eta\Lambda=\eta
  \qquad\text{(Lorentzian)}.
\]
Finally,
\[
  \bar R(M):=R(\Lambda),
\]
where \(R(\Lambda)\) is the representation matrix associated with the
\(SO(D)\) rotation or Lorentz transformation \(\Lambda^\mu{}_\nu\):
\[
  R(\Lambda)
  =
  \exp\left(-\frac{\ii}{2}\omega_{\mu\nu}S^{\mu\nu}\right),
  \qquad
  \Lambda^\mu{}_\nu=\delta^\mu{}_\nu+\omega^\mu{}_\nu .
\]
This is the general conformal transformation property of a primary.
Conformal transformations of descendants can be deduced from the
transformation of the primary together with repeated use of the
conformal-algebra commutation relations.

As a special case, if the primaries transform in a tensor
representation of the Lorentz group, for example
\(\phi_{\mu_1\cdots\mu_s}(x)\), the conformal transformation can be
written explicitly as
\[
  \widetilde\phi_{\mu_1\cdots\mu_s}(\widetilde x)
  =
  \left(
    \det\frac{\partial\widetilde x}{\partial x}
  \right)^{-(\Delta-s)/D}
  \frac{\partial x^{\nu_1}}{\partial\widetilde x^{\mu_1}}
  \cdots
  \frac{\partial x^{\nu_s}}{\partial\widetilde x^{\mu_s}}
  \phi_{\nu_1\cdots\nu_s}(x).
\]
Equivalently,
\[
  \widetilde\phi^{\mu_1\cdots\mu_s}(\widetilde x)
  =
  \left(
    \det\frac{\partial\widetilde x}{\partial x}
  \right)^{-(\Delta+s)/D}
  \frac{\partial\widetilde x^{\mu_1}}{\partial x^{\nu_1}}
  \cdots
  \frac{\partial\widetilde x^{\mu_s}}{\partial x^{\nu_s}}
  \phi^{\nu_1\cdots\nu_s}(x).
\]

\subsection{Consequences of Unitarity}

Under the state/operator map,
\[
  \widehat{\mathcal O}(0)
  \longleftrightarrow
  \ket{\mathcal O}\in\mathcal H_{S^{d-1}}.
\]
Unitarity implies that, for any \(a^\mu\),
\[
  \norm{a^\mu\widehat P_\mu\ket{\mathcal O}}^2\ge 0.
\]
Equivalently, the Hermitian matrix
\[
  \bra{\mathcal O}\widehat P_\mu^\dagger\widehat P_\nu
  \ket{\mathcal O}
\]
is positive semidefinite.  Using
\[
  \widehat P_\mu^\dagger=\widehat K_\mu,
\]
this matrix is
\[
  \bra{\mathcal O}\widehat K_\mu\widehat P_\nu
  \ket{\mathcal O}.
\]

Now suppose \(\mathcal O=\phi\) is a primary.  Since
\(\widehat K_\mu\ket{\phi}=0\), it follows that
\[
\begin{aligned}
  0
  &\le
  \bra{\phi}\widehat K_\mu\widehat P_\nu\ket{\phi}  \\
  &=
  \bra{\phi}[\widehat K_\mu,\widehat P_\nu]\ket{\phi}  \\
  &=
  -2\ii
  \bra{\phi}
  \left(\delta_{\mu\nu}\widehat D-\widehat J_{\mu\nu}\right)
  \ket{\phi}.
\end{aligned}
\]
Thus
\[
  \Delta_\phi
  \ge
  \text{any eigenvalue of }
  \frac{\ii\bra{\phi}\widehat J_{\mu\nu}\ket{\phi}}
       {\braket{\phi|\phi}}.
\]
For example, if \(\phi\) is a scalar primary, then
\[
  \widehat J_{\mu\nu}\ket{\phi}=0,
  \qquad
  \Delta_\phi\ge 0.
\]

There are more constraints.  Consider
\[
  M^{(2)}_{\alpha\beta,\mu\nu}
  \equiv
  \bra{\phi}
  \widehat K_\alpha\widehat K_\beta
  \widehat P_\mu\widehat P_\nu
  \ket{\phi}.
\]
The matrix \(M^{(2)}\) is also positive semidefinite.  Repeatedly using
\([\widehat K,\widehat P]=\cdots\) and
\(\widehat K_\mu\ket{\phi}=0\), one finds for a scalar primary
\[
  M^{(2)}_{\alpha\beta,\mu\nu}
  =
  4\left[
    \left(
      \delta_{\alpha\mu}\delta_{\beta\nu}
      +
      \delta_{\alpha\nu}\delta_{\beta\mu}
    \right)\Delta_\phi(\Delta_\phi+1)
    -
    \delta_{\alpha\beta}\delta_{\mu\nu}\Delta_\phi
  \right]\braket{\phi|\phi}.
\]
Contracting both sides with
\(\delta^{\alpha\beta}\delta^{\mu\nu}\), and writing the spacetime
dimension as \(d\) to avoid confusion with the dilatation generator, gives
\[
\begin{aligned}
  0
  &\le
  \delta^{\alpha\beta}\delta^{\mu\nu}
  M^{(2)}_{\alpha\beta,\mu\nu}  \\
  &=
  4\left[2d\Delta_\phi(\Delta_\phi+1)-d^2\Delta_\phi\right]
  \braket{\phi|\phi}.
\end{aligned}
\]
There are therefore two possibilities:
\[
  \Delta_\phi=0
  \quad\Longrightarrow\quad
  \widehat P_\mu\ket{\phi}
  =
  \ii\ket{\partial_\mu\phi}=0,
\]
so that
\[
  \partial_\mu\phi(x)=0,
  \qquad
  \phi(x)\propto 1,
\]
the identity operator; or
\[
  \Delta_\phi\ge\frac{d-2}{2}.
\]
This is the unitarity bound for a scalar primary.

The scalar unitarity bound is saturated,
\[
  \Delta_\phi=\frac{d-2}{2},
\]
if and only if
\[
  \delta^{\alpha\beta}\delta^{\mu\nu}
  M^{(2)}_{\alpha\beta,\mu\nu}=0.
\]
Equivalently,
\[
  \widehat P^\mu\widehat P_\mu\ket{\phi}=0,
\]
or, in local-operator language,
\[
  \partial^\mu\partial_\mu\phi(x)=0.
\]
This is the massless Klein--Gordon equation, so \(\phi\) is a free
massless scalar field.

For non-scalar primaries, for example primaries in a rank-\(s\)
symmetric traceless tensor representation of the \(SO(d)\) rotation
group,
\[
  \phi_{\mu_1\cdots\mu_s}(x),
\]
with
\[
  \phi_{\mu_1\cdots\mu_s}
  \text{ symmetric in }(\mu_1,\ldots,\mu_s),
  \qquad
  \delta^{\mu_1\mu_2}\phi_{\mu_1\mu_2\cdots\mu_s}=0,
\]
the spin-\(s\) unitarity bound is
\[
  \Delta_\phi\ge s+d-2.
\]

\section{Correlation Functions and Conformal Frames}

The Ward identities strongly constrain correlation functions.  If
\(\mathcal O_i\) are conformal primary operators, then under a conformal
transformation \(x\mapsto x'\) their correlator transforms by the product
of the local scale factors,
\[
  \left\langle
    \mathcal O_1(x_1)\cdots\mathcal O_n(x_n)
  \right\rangle
  =
  \prod_{i=1}^n
  \Omega(x_i')^{-\Delta_i}
  \left\langle
    \mathcal O_1(x_1')\cdots\mathcal O_n(x_n')
  \right\rangle .
\]
Equivalently, the conformal Ward identity reads
\[
  \sum_{i=1}^n
  \left\langle
    \mathcal O_1(x_1)\cdots
    \delta_\epsilon\mathcal O_i(x_i)\cdots
    \mathcal O_n(x_n)
  \right\rangle=0.
\]

For a scalar primary \(\phi(x)\),
\[
  \phi(x)\longmapsto \phi'(x')
  =
  \left|\det\frac{\partial x'}{\partial x}\right|^{-\Delta_\phi/d}
  \phi(x).
\]
The one-point function is therefore fixed immediately.  Translation
invariance gives \(\langle\phi(x)\rangle=\langle\phi\rangle\), while
dilatations give
\[
  \langle\phi\rangle
  =
  \lambda^{-\Delta_\phi}\langle\phi\rangle.
\]
Thus
\[
  \langle\phi\rangle=0
  \qquad\text{unless}\qquad
  \Delta_\phi=0,
\]
and in a unitary CFT the \(\Delta=0\) scalar primary is the identity
operator.

\begin{figure}[htbp]
\centering
\begin{tikzpicture}[scale=0.9,>=Stealth]
  % Source family: 253c pp. 52, 54.
  \draw[thick] (-5.5,-1.75) rectangle (-2.05,1.65);
  \node[anchor=west] at (-5.45,1.95) {\(\mathbb R^d\)};
  \foreach \x/\y/\lab in {-4.85/0.45/{x_1},-3.7/-0.55/{x_2},-2.75/0.8/{x_n}} {
    \node[qft insertion] at (\x,\y) {};
    \node[above right,violet] at (\x,\y) {\(\mathcal O(\lab)\)};
  }
  \draw[qft surface] (-4.85,0.45) ellipse (0.62 and 0.36);
  \node[green!45!black] at (-5.05,1.15) {\(\Sigma_1\)};
  \draw[qft contour] (-3.5,-0.2) .. controls (-4.2,-1.4) and (-2.6,-1.35) .. (-2.35,-0.25)
    .. controls (-2.3,0.8) and (-3.2,1.25) .. (-3.5,-0.2);
  \node[violet] at (-2.55,-1.25) {\(U\)};

  \draw[qft arrow] (-1.45,0) -- (-0.2,0) node[midway,above] {\(x\mapsto x'\)};

  \draw[thick] (0.75,-1.75) rectangle (4.25,1.65);
  \node[anchor=west] at (0.8,1.95) {\(\mathbb R^d\)};
  \foreach \x/\y/\lab in {1.35/-0.7/{x'_1},2.55/0.8/{x'_2},3.65/-0.05/{x'_n}} {
    \node[qft insertion] at (\x,\y) {};
    \node[above right,violet] at (\x,\y) {\(\mathcal O(\lab)\)};
  }
  \draw[qft surface] (1.35,-0.7) ellipse (0.46 and 0.28);
  \node[green!45!black] at (0.95,-0.12) {\(\Sigma_1'\)};
  \draw[qft contour] (2.05,-1.15) .. controls (1.6,0.2) and (2.8,1.35) .. (3.85,0.95)
    .. controls (4.0,-0.55) and (3.05,-1.45) .. (2.05,-1.15);
  \node[violet] at (3.8,1.28) {\(U'\)};
  \node[qft note,blue!70!black] at (0.15,-2.05)
    {conformal maps move insertions and rescale each primary};
\end{tikzpicture}
\caption{Conformal covariance of local insertions and the Ward-surface picture.}
\end{figure}

For scalar primaries, translations, rotations, dilatations, and special
conformal transformations fix the two-point function up to an overall
normalization:
\[
  \left\langle\phi_1(x)\phi_2(y)\right\rangle
  =
  \frac{N_{12}}{|x-y|^{\Delta_1+\Delta_2}},
  \qquad
  N_{12}=0\quad\text{unless}\quad \Delta_1=\Delta_2 .
\]
When \(\Delta_1=\Delta_2=\Delta\), this becomes
\[
  \left\langle\phi_1(x)\phi_2(y)\right\rangle
  =
  \frac{N_{12}}{|x-y|^{2\Delta}}.
\]
One way to see the equality of dimensions is to use inversion,
\[
  x^\mu\mapsto x'^\mu=\frac{x^\mu}{x^2},
  \qquad
  \det\frac{\partial x'}{\partial x}=-(x^2)^{-d},
\]
where the parity sign is irrelevant for parity-even scalar primaries.
Under inversion,
\[
  |x'-y'|^2=\frac{|x-y|^2}{x^2y^2}.
\]
Invariance of the two-point function therefore gives
\[
  N_{12}\,
  \frac{|x|^{-2\Delta_1}|y|^{-2\Delta_2}}
       {|x'-y'|^{\Delta_1+\Delta_2}}
  =
  \frac{N_{12}}{|x-y|^{\Delta_1+\Delta_2}},
\]
which is possible for nonzero \(N_{12}\) only when
\(\Delta_1=\Delta_2\).

Under the state/operator map, \(x=r\Omega\), \(r=\ee^\tau\), and
\[
  \phi(x)=r^{-\Delta_\phi}\widetilde\phi(\tau,\Omega).
\]
Consequently
\[
  \left\langle
    \widetilde\phi_1(\tau_1,\Omega_1)
    \widetilde\phi_2(\tau_2,\Omega_2)
  \right\rangle
  =
  \frac{N_{12}}
  {\left(2\cosh(\tau_{12})-2\Omega_1\cdot\Omega_2\right)^\Delta}
  \quad(\Delta_1=\Delta_2=\Delta),
\]
which is manifestly invariant under translations of cylinder time.
The BPZ conjugate bra is obtained by sending the insertion to
\(\tau\to+\infty\):
\[
  \bra{\phi}
  =
  \lim_{\tau\to+\infty}
  \ee^{\Delta\tau}\bra{\Omega}\,\widetilde\phi(\tau,\Omega),
  \qquad
  \ket{\phi}
  =
  \lim_{\tau\to-\infty}
  \ee^{-\Delta\tau}\widetilde\phi(\tau,\Omega)\ket{\Omega}.
\]
Thus \(N_{12}\) is the BPZ inner product
\[
  N_{12}=\braket{\phi_1|\phi_2},
\]
and one may choose a Hermitian basis of primaries with
\(\braket{\phi_i|\phi_j}=\delta_{ij}\).

\begin{figure}[htbp]
\centering
\begin{tikzpicture}[scale=0.9,>=Stealth]
  % Source family: 253c pp. 58--60.
  \draw[thick] (-5.5,-1.45) rectangle (-2.3,1.55);
  \node[anchor=west] at (-5.45,1.82) {\(\mathbb R^d\)};
  \draw[qft surface] (-3.9,0.0) circle (0.8);
  \node[qft insertion] at (-4.18,-0.1) {};
  \node[violet] at (-4.48,0.22) {\(\phi(0)\)};
  \node[qft insertion] at (-3.36,0.46) {};
  \node[violet] at (-2.9,0.65) {\(\phi(x)\)};
  \draw[qft flow] (-3.9,0) -- (-3.22,0.52) node[midway,above] {\(r\)};

  \draw[qft arrow] (-1.55,0) -- (-0.35,0) node[midway,above] {\(r=\ee^\tau\)};

  \draw[thick] (0.55,1.35) ellipse (0.85 and 0.22);
  \draw[thick] (-0.3,1.35) -- (-0.3,-1.35);
  \draw[thick] (1.4,1.35) -- (1.4,-1.35);
  \draw[thick] (0.55,-1.35) ellipse (0.85 and 0.22);
  \draw[qft surface] (-0.3,0.0) arc[start angle=180,end angle=360,x radius=0.85,y radius=0.22];
  \draw[qft surface,dashed] (-0.3,0.0) arc[start angle=180,end angle=0,x radius=0.85,y radius=0.22];
  \node[violet] at (0.55,-0.45) {\(\ket{\phi}\)};
  \draw[qft arrow] (1.8,0) -- (3.0,0) node[midway,above] {BPZ};
  \draw[thick] (4.05,1.35) ellipse (0.85 and 0.22);
  \draw[thick] (3.2,1.35) -- (3.2,-1.35);
  \draw[thick] (4.9,1.35) -- (4.9,-1.35);
  \draw[thick] (4.05,-1.35) ellipse (0.85 and 0.22);
  \node[violet] at (4.05,0.55) {\(\bra{\phi_1}\)};
  \node[violet] at (4.05,-0.62) {\(\ket{\phi_2}\)};
  \draw[qft contour] (3.25,-0.05) arc[start angle=185,end angle=355,x radius=0.8,y radius=0.22];
  \draw[qft contour,dashed] (3.25,-0.05) arc[start angle=185,end angle=5,x radius=0.8,y radius=0.22];
  \node[qft note,blue!70!black] at (0.9,-1.88)
    {the two-point coefficient is the BPZ inner product};
\end{tikzpicture}
\caption{State/operator map and BPZ conjugation for two-point functions.}
\end{figure}

For non-scalar primaries, the same analysis determines the tensor
structure.  A useful example is the stress tensor:
\[
  \left\langle T_{\mu\nu}(x)T_{\rho\sigma}(0)\right\rangle
  =
  \frac{C_T}{|x|^{2d}}\,I_{\mu\nu,\rho\sigma}(x),
\]
where
\[
  I_{\mu\nu,\rho\sigma}(x)
  =
  \frac12\left(I_{\mu\rho}(x)I_{\nu\sigma}(x)
  +I_{\mu\sigma}(x)I_{\nu\rho}(x)\right)
  -\frac1d\delta_{\mu\nu}\delta_{\rho\sigma},
  \qquad
  I_{\mu\nu}(x)=\delta_{\mu\nu}-2\frac{x_\mu x_\nu}{x^2}.
\]
Here \(T_{\mu\nu}\) is symmetric, traceless, and conserved, so its
dimension is \(\Delta_T=d\).  The coefficient \(C_T\) is intrinsic CFT
data because the normalization of \(T_{\mu\nu}\) is fixed by the Ward
identity for translations.

Three-point functions are similarly constrained.  For the stress tensor
itself,
\[
  \left\langle
    T_{\mu\nu}(x_1)T_{\rho\sigma}(x_2)T_{\alpha\beta}(x_3)
  \right\rangle
  =
  \sum_{i=1}^3
  a_i\,\mathcal I^{(i)}_{\mu\nu,\rho\sigma,\alpha\beta}(x_{12},x_{13}),
\]
where the \(\mathcal I^{(i)}\) are the three known conformal tensor
structures.  Conservation fixes one linear combination of the \(a_i\)'s
in terms of \(C_T\), since \(T\) appearing in the \(TT\) OPE is a
descendant determined by the conformal algebra.

For two scalar primaries and the stress tensor, conformal symmetry and
current conservation give
\[
\begin{aligned}
  &\left\langle
    \phi_1(x_1)\phi_2(x_2)T_{\mu\nu}(x_3)
  \right\rangle  \\
  &\qquad =
  \frac{C_{12T}}
  {|x_{13}|^d|x_{23}|^d|x_{12}|^{2\Delta-d}}
  \left(
    \frac{X_\mu X_\nu}{X^2}-\frac{\delta_{\mu\nu}}d
  \right),
\end{aligned}
\]
where
\[
  X^\mu=\frac{x_{13}^\mu}{x_{13}^2}
       -\frac{x_{23}^\mu}{x_{23}^2},
  \qquad
  \Delta_{\phi_1}=\Delta_{\phi_2}\equiv\Delta.
\]
It vanishes when \(\Delta_{\phi_1}\ne\Delta_{\phi_2}\).  The Ward
identity for dilatations determines its coefficient in terms of the
two-point normalization:
\[
  C_{12T}
  =
  -\frac{\Delta\,N_{12}}
  {\left(1-\frac1d\right)A_{d-1}},
  \qquad
  A_{d-1}=\frac{2\pi^{d/2}}{\Gamma(d/2)}.
\]

\begin{figure}[htbp]
\centering
\begin{tikzpicture}[scale=0.88,>=Stealth]
  % Source family: 253c pp. 64--66.
  \node at (-3.9,1.75) {\(\langle T_{\mu\nu}(x_1)T_{\rho\sigma}(x_2)T_{\alpha\beta}(x_3)\rangle\)};
  \foreach \x/\y/\lab in {-4.9/0/{x_1},-2.9/0/{x_2},-3.9/-1.35/{x_3}} {
    \node[qft insertion] at (\x,\y) {};
    \node[above,violet] at (\x,\y) {\(\lab\)};
  }
  \draw[qft contour] (-4.9,0) circle (0.62);
  \draw[qft contour] (-2.9,0) circle (0.62);
  \draw[qft contour] (-3.9,-1.35) circle (0.62);
  \node[qft note,green!45!black] at (-3.9,-2.25) {three independent tensor structures};

  \draw[qft arrow] (-0.9,-0.25) -- (0.65,-0.25) node[midway,above] {Ward identity};

  \draw[thick] (2.4,0.0) circle (1.05);
  \draw[qft contour] (2.4,0.0) circle (0.72);
  \node[qft insertion] at (2.4,0.0) {};
  \node[violet] at (2.82,0.22) {\(\phi\)};
  \draw[qft flow] (1.25,1.1) .. controls (1.7,1.6) and (2.65,1.65) .. (3.2,1.12);
  \node[blue!70!black] at (2.25,1.78) {\(T_{\mu\nu}\)};
  \node at (2.4,-1.55) {\(\partial^\mu T_{\mu\nu}=0\)};
  \node[qft note,orange!85!black] at (2.4,-2.15)
    {conservation fixes scalar--scalar--\(T\) normalization};
\end{tikzpicture}
\caption{Stress-tensor correlators are fixed up to a small number of CFT data.}
\end{figure}

The three-point function of scalar primaries is fixed up to the OPE
coefficient \(C_{123}\):
\[
  \left\langle \phi_1(x_1)\phi_2(x_2)\phi_3(x_3)\right\rangle
  =
  \frac{C_{123}}
  {|x_{12}|^{\Delta_1+\Delta_2-\Delta_3}
   |x_{23}|^{\Delta_2+\Delta_3-\Delta_1}
   |x_{13}|^{\Delta_1+\Delta_3-\Delta_2}}.
\]
For four points, conformal symmetry leaves two independent cross ratios,
\[
  u=\frac{x_{12}^2x_{34}^2}{x_{13}^2x_{24}^2},
  \qquad
  v=\frac{x_{14}^2x_{23}^2}{x_{13}^2x_{24}^2}.
\]
Equivalently, after a conformal map the four points can be placed in a
two-plane and then at
\[
  z_1=0,\qquad z_2=z,\qquad z_3=1,\qquad z_4=\infty .
\]
The complex cross ratio is
\[
  z=\frac{z_{12}z_{34}}{z_{13}z_{24}},
  \qquad
  u=z\bar z,\qquad v=(1-z)(1-\bar z).
\]
The scalar four-point function is therefore fixed up to one arbitrary
function of the cross ratios:
\[
\begin{aligned}
  &\left\langle
    \phi_1(x_1)\phi_2(x_2)\phi_3(x_3)\phi_4(x_4)
  \right\rangle  \\
  &\quad =
  \frac{g(u,v)}
  {|x_{12}|^{\Delta_1+\Delta_2}
   |x_{34}|^{\Delta_3+\Delta_4}}
  \left(\frac{|x_{24}|}{|x_{14}|}\right)^{\Delta_1-\Delta_2}
  \left(\frac{|x_{14}|}{|x_{13}|}\right)^{\Delta_3-\Delta_4}.
\end{aligned}
\]

\begin{figure}[htbp]
\centering
\begin{tikzpicture}[scale=0.88,>=Stealth]
  % Source family: 253c pp. 67--69.
  \draw[thick] (-5.2,-1.45) rectangle (-2.0,1.55);
  \foreach \x/\y/\lab in {-4.75/-0.75/{x_1},-4.25/0.8/{x_2},-3.1/0.45/{x_3},-2.55/-0.65/{x_4}} {
    \node[qft insertion] at (\x,\y) {};
    \node[above right,violet] at (\x,\y) {\(\lab\)};
  }
  \draw[qft arrow] (-1.45,0) -- (-0.1,0) node[midway,above] {conformal map};
  \draw[thick] (0.65,-1.25) -- (0.95,1.25) -- (3.65,1.1) -- (3.35,-1.35) -- cycle;
  \foreach \x/\y/\lab in {1.0/-0.75/{0},1.25/0.7/{z},2.95/0.6/{1},3.1/-0.7/{\infty}} {
    \node[qft insertion] at (\x,\y) {};
    \node[above right,violet] at (\x,\y) {\(\lab\)};
  }
  \node[qft note,blue!70!black] at (1.9,-1.85)
    {four points can be placed at \(0,z,1,\infty\)};
  \node[orange!85!black] at (2.1,1.65) {\(z,\bar z\)};
\end{tikzpicture}
\caption{Conformal frame for scalar four-point functions.}
\end{figure}

\section{Operator Product Expansion and Conformal Blocks}

The operator product expansion is the statement that products of local
operators can be expanded in a basis of local operators:
\[
  \phi_1(x)\phi_2(0)
  =
  \sum_{\mathcal O}
  C_{12\mathcal O}(x,\partial)\,\mathcal O(0).
\]
The sum runs over primary operators and their descendants, organized by
increasing scaling dimension.
For scalar \(\phi_1,\phi_2\), a spin-\(\ell\) symmetric traceless
primary \(\mathcal O_{\mu_1\cdots\mu_\ell}\) can appear with leading
coefficient
\[
  \phi_1(x)\phi_2(0)
  \supset
  \lambda_{12\mathcal O}\,
  \frac{x^{\mu_1}\cdots x^{\mu_\ell}}
       {|x|^{\Delta_1+\Delta_2-\Delta+\ell}}\,
  \mathcal O_{\mu_1\cdots\mu_\ell}(0)
  +\text{descendants}.
\]
Only symmetric traceless tensors can appear in this scalar-scalar OPE:
antisymmetric tensors vanish when contracted with powers of the same
vector \(x^\mu\), and traces are removed.

The coefficients of descendants are not independent.  They are fixed by
conformal symmetry from the three-point function
\(\langle\phi_1(x)\phi_2(0)\mathcal O(y)\rangle\).  For example, for a
scalar primary \(\mathcal O\),
\[
  \phi_1(x)\phi_2(0)
  \supset
  \lambda_{12\mathcal O}
  |x|^{\Delta-\Delta_1-\Delta_2}
  \left(
    1+
    \frac{\Delta+\Delta_1-\Delta_2}{2\Delta}
    x^\mu\partial_\mu+\cdots
  \right)\mathcal O(0).
\]
Thus the OPE data are the primary spectrum and the structure constants;
descendant terms are representation-theoretic consequences of the
conformal algebra.

A primary, together with all of its descendants, spans an irreducible
representation \(V_\mathcal O\) of the conformal algebra.  The OPE is the
decomposition of the product representation generated by
\(\phi_1(x_1)\phi_2(x_2)\) into such irreducibles:
\[
  V_{\phi_1}\otimes V_{\phi_2}
  =
  \bigoplus_{\mathcal O} m_{\mathcal O} V_{\mathcal O}.
\]
The multiplicity \(m_{\mathcal O}\) is the number of independent tensor
structures allowed in the corresponding three-point function.

\begin{figure}[htbp]
\centering
\begin{tikzpicture}[scale=0.9,>=Stealth]
  % Source family: 253c pp. 69--73.
  \draw[thick] (-5.25,1.3) ellipse (0.82 and 0.22);
  \draw[thick] (-6.07,1.3) -- (-6.07,-1.3);
  \draw[thick] (-4.43,1.3) -- (-4.43,-1.3);
  \draw[thick] (-5.25,-1.3) ellipse (0.82 and 0.22);
  \draw[qft contour] (-6.0,0.2) arc[start angle=185,end angle=355,x radius=0.75,y radius=0.2];
  \draw[qft contour,dashed] (-6.0,0.2) arc[start angle=185,end angle=5,x radius=0.75,y radius=0.2];
  \node[qft insertion] at (-5.52,0.15) {};
  \node[qft insertion] at (-5.0,0.35) {};
  \node[violet] at (-5.55,0.62) {\(\phi_1\)};
  \node[violet] at (-4.72,0.62) {\(\phi_2\)};
  \node[orange!85!black] at (-5.25,-0.55) {local expansion};

  \draw[qft arrow] (-3.65,0) -- (-2.35,0) node[midway,above] {OPE};

  \draw[thick] (-1.25,1.3) ellipse (0.82 and 0.22);
  \draw[thick] (-2.07,1.3) -- (-2.07,-1.3);
  \draw[thick] (-0.43,1.3) -- (-0.43,-1.3);
  \draw[thick] (-1.25,-1.3) ellipse (0.82 and 0.22);
  \foreach \y/\lab in {0.72/{\mathcal O},0.25/{\partial\mathcal O},-0.22/{\partial^2\mathcal O}} {
    \node[qft insertion] at (-1.25,\y) {};
    \node[violet,right] at (-1.05,\y) {\(\lab\)};
  }
  \draw[qft arrow] (0.35,0) -- (1.65,0) node[midway,above] {basis};
  \node at (3.25,0.65) {\(\displaystyle
    \phi_1\phi_2=\sum_{\mathcal O}\lambda_{12\mathcal O}
    \bigl(\mathcal O+\partial\mathcal O+\cdots\bigr)\)};
  \node[qft note,blue!70!black] at (2.95,-0.95)
    {descendants are fixed by conformal symmetry};
\end{tikzpicture}
\caption{The OPE replaces nearby insertions by primaries and descendants.}
\end{figure}

Applying the OPE inside a four-point function gives
\[
  \left\langle\phi_1\phi_2\phi_3\phi_4\right\rangle
  =
  \sum_{\mathcal O}
  \lambda_{12\mathcal O}\lambda_{34\mathcal O}
  G_{\Delta,\ell}(u,v),
\]
where \(G_{\Delta,\ell}\) is a conformal block.
With an orthonormal basis of primaries, \(N_{ij}=\delta_{ij}\), the CFT
data
\[
  \{\Delta_i,\ \ell_i,\ \lambda_{ijk}\}
\]
determine all Euclidean Green functions by repeated use of the OPE.  The
key consistency condition is associativity of the OPE, equivalently
crossing symmetry of four-point functions.

In radial quantization, write
\[
  z=r\ee^{\ii\theta},\qquad
  \bar z=r\ee^{-\ii\theta}.
\]
For identical external scalars, the conformal block expansion has the
form
\[
  g(u,v)=1+\sum_{\mathcal O\ne\mathbf1}
  \lambda_{\mathcal O}^2\,G_{\Delta,\ell}(u,v),
  \qquad
  \lambda_{\mathcal O}^2\ge0.
\]
The identity contribution is the leading \(1\).  At each descendant
level \(n\), if \(\ket{\mathcal O;N}\) is a basis of level-\(n\)
descendants and
\[
  G_{NM}=\braket{\mathcal O;N|\mathcal O;M}
\]
is the Gram matrix, then the contribution of this primary family is
schematically
\[
  \sum_{N,M}
  \bra{\Omega}\phi\phi\ket{\mathcal O;N}
  (G^{-1})_{NM}
  \bra{\mathcal O;M}\phi\phi\ket{\Omega}.
\]
If the Gram matrix degenerates, one first passes to a reduced linearly
independent basis.

The leading small-\(r\) behavior is fixed by the primary state:
\[
  G_{\Delta,\ell}(r,\eta)
  =
  r^\Delta C_\ell^{(d/2-1)}(\eta)
  +\text{higher powers of }r,
  \qquad
  \eta=\cos\theta,
\]
where \(C_\ell^{(d/2-1)}\) is a Gegenbauer polynomial.  In special
dimensions this reduces to familiar polynomials: in \(d=2\) to
\(\cos(\ell\theta)\), in \(d=3\) to Legendre polynomials, and in \(d=4\)
to \(\sin((\ell+1)\theta)/\sin\theta\).

\begin{figure}[htbp]
\centering
\begin{tikzpicture}[scale=0.88,>=Stealth]
  % Source family: 253c pp. 76--80, 85, 97.
  \draw[thick] (-5.6,0) circle (1.15);
  \node[qft insertion] at (-6.2,0.35) {};
  \node[qft insertion] at (-5.85,-0.4) {};
  \node[qft insertion] at (-4.9,0.45) {};
  \node[qft insertion] at (-4.62,-0.45) {};
  \draw[qft contour] (-6.03,0.0) circle (0.62);
  \draw[qft contour] (-4.76,0.0) circle (0.62);
  \node[blue!70!black] at (-6.65,0.95) {OPE};
  \node[blue!70!black] at (-4.2,0.95) {OPE};
  \node at (-5.6,-1.55) {\(s\)-channel};

  \draw[qft arrow] (-3.95,0) -- (-2.9,0);

  \draw[thick] (-1.7,0) circle (1.15);
  \node[qft insertion] at (-2.25,0.35) {};
  \node[qft insertion] at (-1.95,-0.4) {};
  \node[qft insertion] at (-0.95,0.45) {};
  \node[qft insertion] at (-0.68,-0.45) {};
  \draw[qft contour] (-2.08,0.0) circle (0.62);
  \draw[qft contour] (-0.82,0.0) circle (0.62);
  \draw[qft flow] (-2.08,0) .. controls (-1.65,0.55) and (-1.18,0.55) .. (-0.82,0);
  \node at (-1.7,-1.55) {\(\sum_{\mathcal O}\lambda_{12\mathcal O}\lambda_{34\mathcal O}G_{\mathcal O}\)};

  \draw[qft arrow] (0.45,0) -- (1.5,0);

  \begin{scope}[shift={(3.2,0)}]
    \node[qft insertion] (a) at (-0.9,0.7) {};
    \node[qft insertion] (b) at (-0.9,-0.7) {};
    \node[qft insertion] (c) at (0.9,0.7) {};
    \node[qft insertion] (d) at (0.9,-0.7) {};
    \draw[qft feynman] (a) -- (0,0);
    \draw[qft feynman] (b) -- (0,0);
    \draw[qft feynman] (0,0) -- (c);
    \draw[qft feynman] (0,0) -- (d);
    \node[orange!85!black] at (0,0.28) {\(\mathcal O\)};
    \node at (0,-1.55) {block as exchange};
  \end{scope}
\end{tikzpicture}
\caption{OPE channel decomposition of a four-point function.}
\end{figure}

The two different ways of grouping the four points lead to crossing
symmetry.  Schematically,
\[
  \sum_{\mathcal O}\lambda_{12\mathcal O}\lambda_{34\mathcal O}
  G_{\Delta,\ell}(u,v)
  =
  \sum_{\mathcal O}\lambda_{14\mathcal O}\lambda_{23\mathcal O}
  G_{\Delta,\ell}(v,u).
\]

\begin{figure}[htbp]
\centering
\begin{tikzpicture}[scale=0.92,>=Stealth]
  % Source family: 253c pp. 78--79, 85.
  \begin{scope}[shift={(-3.8,0)}]
    \node[qft insertion] (a) at (-1,0.8) {};
    \node[qft insertion] (b) at (-1,-0.8) {};
    \node[qft insertion] (c) at (1,0.8) {};
    \node[qft insertion] (d) at (1,-0.8) {};
    \draw[qft feynman] (a) -- (0,0) -- (c);
    \draw[qft feynman] (b) -- (0,0) -- (d);
    \node[blue!70!black] at (-1.2,1.15) {\(\phi_1\)};
    \node[blue!70!black] at (-1.2,-1.15) {\(\phi_2\)};
    \node[blue!70!black] at (1.2,1.15) {\(\phi_3\)};
    \node[blue!70!black] at (1.2,-1.15) {\(\phi_4\)};
    \node at (0,-1.55) {\(s\)-channel};
  \end{scope}
  \node at (0,0) {\(=\)};
  \begin{scope}[shift={(3.0,0)}]
    \node[qft insertion] (a) at (-1,0.8) {};
    \node[qft insertion] (b) at (-1,-0.8) {};
    \node[qft insertion] (c) at (1,0.8) {};
    \node[qft insertion] (d) at (1,-0.8) {};
    \draw[qft feynman] (a) -- (0,0.55) -- (d);
    \draw[qft feynman] (b) -- (0,-0.55) -- (c);
    \draw[qft feynman] (0,0.55) -- (0,-0.55);
    \node[blue!70!black] at (-1.2,1.15) {\(\phi_1\)};
    \node[blue!70!black] at (-1.2,-1.15) {\(\phi_2\)};
    \node[blue!70!black] at (1.2,1.15) {\(\phi_3\)};
    \node[blue!70!black] at (1.2,-1.15) {\(\phi_4\)};
    \node at (0,-1.55) {\(t\)-channel};
  \end{scope}
  \node[qft note,orange!85!black] at (0,1.85)
    {crossing equates different OPE decompositions};
\end{tikzpicture}
\caption{Crossing symmetry as equality of OPE channels.}
\end{figure}

In radial quantization, it is often useful to use the variable
\[
  \rho=\frac{z}{(1+\sqrt{1-z})^2},
\]
which maps the cut \(z\)-plane into the unit disk and improves the
convergence of the OPE.
The inverse relation is
\[
  z=\frac{4\rho}{(1+\rho)^2}.
\]
The crossing-symmetric point \(z=\bar z=1/2\) is mapped to
\[
  \rho=3-2\sqrt2\simeq0.17,
\]
so the radial expansion is rapidly convergent there.

For identical scalar primaries, crossing gives
\[
  v^{\Delta_\phi}g(u,v)=u^{\Delta_\phi}g(v,u).
\]
Equivalently,
\[
  v^{\Delta_\phi}-u^{\Delta_\phi}
  +
  \sum_{\mathcal O\ne\mathbf1}
  \lambda_{\mathcal O}^2
  \left(
    v^{\Delta_\phi}G_{\Delta,\ell}(u,v)
    -u^{\Delta_\phi}G_{\Delta,\ell}(v,u)
  \right)=0.
\]
If one assumes a large gap above the identity in the \(\phi\times\phi\)
OPE, expanding near \(z=\bar z=1/2\) leads to a contradiction with
positivity of \(\lambda_{\mathcal O}^2\).  Hence there must be an upper
bound on the first nontrivial operator dimension as a function of
\(\Delta_\phi\).  Making this argument sharp requires accurate
conformal blocks.

\begin{figure}[htbp]
\centering
\begin{tikzpicture}[scale=0.9,>=Stealth]
  % Source family: 253c pp. 79--83.
  \draw[thick] (-5.5,-1.2) rectangle (-2.6,1.2);
  \draw[qft surface] (-4.05,0) circle (0.75);
  \draw[dashed,blue!70!black] (-4.05,0) circle (0.42);
  \node[qft insertion] at (-4.42,0.0) {};
  \node[qft insertion] at (-3.63,0.0) {};
  \node[violet] at (-4.7,0.45) {\(x_1,x_2\)};
  \node[violet] at (-3.22,0.45) {\(x_3,x_4\)};
  \node at (-4.05,-1.55) {\(\rho\)-disk};

  \draw[qft arrow] (-1.7,0) -- (-0.35,0) node[midway,above] {cylinder};

  \draw[thick] (0.75,1.25) ellipse (0.8 and 0.2);
  \draw[thick] (-0.05,1.25) -- (-0.05,-1.25);
  \draw[thick] (1.55,1.25) -- (1.55,-1.25);
  \draw[thick] (0.75,-1.25) ellipse (0.8 and 0.2);
  \draw[qft surface] (-0.05,0.35) arc[start angle=180,end angle=360,x radius=0.8,y radius=0.2];
  \draw[qft surface,dashed] (-0.05,0.35) arc[start angle=180,end angle=0,x radius=0.8,y radius=0.2];
  \draw[qft surface] (-0.05,-0.45) arc[start angle=180,end angle=360,x radius=0.8,y radius=0.2];
  \draw[qft surface,dashed] (-0.05,-0.45) arc[start angle=180,end angle=0,x radius=0.8,y radius=0.2];
  \node[violet] at (0.75,0.65) {\(\ket{\Psi}\)};
  \node[violet] at (0.75,-0.75) {\(\bra{\Psi}\)};

  \draw[qft arrow] (2.25,0) -- (3.35,0);
  \node at (4.7,0.28) {\(\displaystyle
    f(u,v)=\sum_{\Delta,\ell}\lambda_{\Delta,\ell}^2
    G_{\Delta,\ell}(u,v)\)};
  \node[qft note,blue!70!black] at (4.65,-0.72)
    {positive norm gives \(\lambda_{\Delta,\ell}^2\ge0\)};
\end{tikzpicture}
\caption{Radial quantization gives a convergent conformal-block expansion.}
\end{figure}

\section{Projective Lightcone and Bootstrap Bounds}

Embedding space realizes the conformal group linearly.  Introduce
coordinates \(X^A\in\mathbb R^{d+1,1}\) satisfying
\[
  X^2=0,\qquad X\sim \lambda X,
\]
so physical points are rays on the projective lightcone.  A scalar
primary is represented by a homogeneous function
\[
  \Phi(\lambda X)=\lambda^{-\Delta}\Phi(X).
\]
A convenient Euclidean section of the cone is
\[
  X^A=(X^+,X^-,X^\mu)=(1,x^2,x^\mu),
\]
with metric chosen so that
\[
  X\cdot Y=-\frac12(x-y)^2.
\]
This immediately reproduces the scalar two- and three-point functions,
since embedding-space correlators can depend only on \(X_i\cdot X_j\)
and the homogeneity weights.

The conformal generators are the \(SO(d+1,1)\) Lorentz generators
\[
  L_{AB}=X_A\frac{\partial}{\partial X^B}
        -X_B\frac{\partial}{\partial X^A},
\]
and the quadratic Casimir
\[
  \mathcal C=-\frac12 L_{AB}L^{AB}
\]
takes the value
\[
  C_{\Delta,\ell}
  =
  \Delta(\Delta-d)+\ell(\ell+d-2)
\]
on the conformal family of a spin-\(\ell\) primary of dimension
\(\Delta\).  Acting on the four-point function, this gives the Casimir
differential equation
\[
  \left[\mathcal D_{z,\bar z}-C_{\Delta,\ell}\right]
  F_{\Delta,\ell}(z,\bar z)=0
\]
for conformal blocks.
More explicitly, the conformal-block contribution in the
\(\phi_1\phi_2\to\phi_3\phi_4\) channel obeys
\[
  \left[-(L_1+L_2)^2-C_{\Delta,\ell}\right]
  G_{\Delta,\ell}(u,v)=0.
\]
In terms of \(z,\bar z\), with \(G_{\Delta,\ell}(u,v)=F_{\Delta,\ell}(z,\bar z)\),
the second-order equation used in the notes is
\[
  \left[
    \mathcal D_z+\mathcal D_{\bar z}
    +(d-2)\frac{z\bar z}{z-\bar z}
    \left((1-z)\partial_z-(1-\bar z)\partial_{\bar z}\right)
  \right]F_{\Delta,\ell}
  =
  \frac12 C_{\Delta,\ell}F_{\Delta,\ell},
\]
where
\[
  \mathcal D_z
  =
  z^2(1-z)\partial_z^2
  +\left(\frac{\Delta_1-\Delta_2-\Delta_3+\Delta_4}{2}-1\right)
  z^2\partial_z
  -\frac{(\Delta_1-\Delta_2)(\Delta_3-\Delta_4)}{4}z,
\]
and \(\mathcal D_{\bar z}\) is the same expression with
\(z\leftrightarrow\bar z\).

In \(d=2\), the global conformal block factorizes into holomorphic and
antiholomorphic hypergeometric functions.  With
\[
  h=\frac{\Delta+\ell}{2},
  \qquad
  \bar h=\frac{\Delta-\ell}{2},
\]
one may write schematically
\[
  F_{\Delta,\ell}(z,\bar z)
  =
  z^h\,
  {}_2F_1\!\left(
    h-\frac{\Delta_1-\Delta_2}{2},
    h+\frac{\Delta_3-\Delta_4}{2};
    2h;
    z
  \right)
  \,
  \bar z^{\bar h}\,
  {}_2F_1\!\left(
    \bar h-\frac{\Delta_1-\Delta_2}{2},
    \bar h+\frac{\Delta_3-\Delta_4}{2};
    2\bar h;
    \bar z
  \right)
  +\text{c.c.},
\]
with the external-dimension differences absorbed into the conventional
definitions.  In \(d=4\) a related hypergeometric expression exists; in
\(d=3\) the scalar blocks are more involved and are usually evaluated
using recursion relations.

For \(d=4\), the notes quote the Dolan--Osborn form.  With
\(h=(\Delta+\ell)/2\) and \(\bar h=(\Delta-\ell)/2\),
\[
\begin{aligned}
  F_{\Delta,\ell}^{(4d)}(z,\bar z)
  &=
  \frac{z^{h+1}\bar z^{\bar h}}{\bar z-z}
  {}_2F_1\!\left(
    h-\frac{\Delta_1-\Delta_2}{2},
    h+\frac{\Delta_3-\Delta_4}{2};
    2h;
    z
  \right) \\
  &\quad\times
  {}_2F_1\!\left(
    \bar h-\frac{\Delta_1-\Delta_2}{2}-1,
    \bar h+\frac{\Delta_3-\Delta_4}{2}-1;
    2\bar h-2;
    \bar z
  \right)
  +(z\leftrightarrow \bar z).
\end{aligned}
\]
The prefactor and the exchanged term make the full expression symmetric
under \(z\leftrightarrow\bar z\).

The notes also record the standard radial-coordinate form of the
three-dimensional recursion.  Write
\[
  G_{\Delta,\ell}(u,v)=(4r)^\Delta h_{\Delta,\ell}(r,\eta),
  \qquad
  \rho=re^{\ii\alpha}
  =
  \frac{z}{\bigl(1+\sqrt{1-z}\bigr)^2},
\]
so that \(\eta=\cos\alpha\).  Then
\[
  h_{\Delta,\ell}
  =
  h_{\infty,\ell}
  +
  \sum_A
  \frac{R_A}{\Delta-\Delta_A^*}
  (4r)^{n_A}h_{\Delta_A^*+n_A,\ell_A}.
\]
Here \(A=(I,n)\), \(I=1,2,3\), \(n\ge1\), with \(n\le\ell\) for
\(I=2\), and
\[
\begin{array}{c|c|c|c}
I & \Delta_A^* & n_A & \ell_A\\ \hline
1 & 1-\ell-n & n & \ell+n\\
2 & \ell+d-1-n & n & \ell-n\\
3 & d/2-n & 2n & \ell
\end{array}
\]
The residues \(R_A\) are products of Pochhammer symbols.  For example,
the first family is written
\[
  R_{1,n}
  =
  -\frac{n(-2)^n}{(n!)^2}
  \left(\frac{\Delta_1-\Delta_2+1-n}{2}\right)_n
  \left(\frac{\Delta_3-\Delta_4+1-n}{2}\right)_n,
\]
where \((x)_n=\Gamma(x+n)/\Gamma(x)\).  The other two residue families
are
\[
\begin{aligned}
  R_{2,n}
  &=
  -\frac{n\,\ell!}{(-2)^n(n!)^2(\ell-n)!}
  \frac{(d+\ell-n-2)_n}
       {\left(\frac d2+\ell-n\right)_n
        \left(\frac d2+\ell-n-1\right)_n} \\
  &\quad\times
  \left(\frac{\Delta_1-\Delta_2+1-n}{2}\right)_n
  \left(\frac{\Delta_3-\Delta_4+1-n}{2}\right)_n ,
  \\
  R_{3,n}
  &=
  -\frac{n(-1)^n\left(\frac d2-n-1\right)_{2n}}
        {(n!)^2
        \left(\frac d2+\ell-n-1\right)_{2n}
        \left(\frac d2+\ell-n\right)_{2n}} \\
  &\quad\times
  \prod_{\sigma=\pm1}
  \left(\frac{\Delta_1-\Delta_2+\sigma(\frac d2+\ell-n)}{2}\right)_n
  \left(\frac{\Delta_3-\Delta_4+\sigma(\frac d2+\ell-n)}{2}\right)_n .
\end{aligned}
\]
The large-\(\Delta\) seed is
\[
\begin{aligned}
  h_{\infty,\ell}(r,\eta)
  &=
  (1-r^2)^{1-d/2}
  (r^2-2r\eta+1)^{-\frac{1-\Delta_1+\Delta_2+\Delta_3-\Delta_4}{2}} \\
  &\qquad\qquad\times
  (r^2+2r\eta+1)^{-\frac{1+\Delta_1-\Delta_2-\Delta_3+\Delta_4}{2}}
  \frac{\ell!}{(-2)^\ell(d/2-1)_\ell}
  C_\ell^{(d/2-1)}(\eta).
\end{aligned}
\]
For identical external scalars
\(\Delta_1=\Delta_2=\Delta_3=\Delta_4=\Delta_\phi\), only even
\(\ell\) appears, some residues vanish, and it suffices to include
\[
  A=(1,2k),\qquad (2,2k),\qquad (3,k),
  \qquad k=1,2,\ldots .
\]

\begin{figure}[htbp]
\centering
\begin{tikzpicture}[scale=0.9,>=Stealth]
  % Source family: 253c pp. 90--94.
  \draw[thick] (-4.7,0) ellipse (1.2 and 0.35);
  \draw[thick] (-5.9,0) -- (-4.7,2.0) -- (-3.5,0);
  \draw[thick] (-5.9,0) -- (-4.7,-2.0) -- (-3.5,0);
  \draw[dashed] (-4.7,0) ellipse (0.55 and 0.16);
  \draw[qft flow] (-4.7,0) -- (-3.62,1.35) node[right] {\(X\)};
  \draw[qft flow] (-4.7,0) -- (-3.95,0.9) node[right] {\(\lambda X\)};
  \node at (-4.7,-2.45) {\(X^2=0,\ X\sim\lambda X\)};

  \draw[qft arrow] (-2.6,0) -- (-1.35,0) node[midway,above] {section};

  \draw[thick] (-0.45,-1.45) -- (2.85,-1.15) -- (3.35,1.25) -- (-0.05,1.0) -- cycle;
  \node[qft insertion] at (0.55,-0.35) {};
  \node[qft insertion] at (1.85,0.45) {};
  \node[violet] at (0.3,-0.72) {\(x\)};
  \node[violet] at (2.1,0.78) {\(y\)};
  \node at (1.45,-1.8) {\(\mathbb R^d\) section};
  \node[qft note,orange!85!black] at (1.35,1.65)
    {\(X\cdot Y\) packages distances};
\end{tikzpicture}
\caption{Projective lightcone formalism for conformal kinematics.}
\end{figure}

The conformal bootstrap studies crossing as an equation for the CFT
data \((\Delta,\ell,\lambda)\).  Positivity of squared OPE coefficients
turns the problem into a linear/semi-definite optimization problem.
A linear functional \(\alpha\) acting on functions of \(z,\bar z\) is
chosen so that
\[
  \alpha(F_{\Delta,\ell})\ge0
\]
for every operator allowed by the assumed spectrum, while
\[
  \alpha(F_{\mathbf1})<0.
\]
Applying \(\alpha\) to the crossing equation then gives a contradiction,
so the assumed spectrum is ruled out.  Numerically, one searches over
finite-dimensional spaces of derivatives at the crossing-symmetric
point.

The continuum of allowed dimensions can be handled by rewriting the
positivity constraints as a polynomial matrix program.  The theorem used
in the notes is that a symmetric polynomial matrix \(M(x)\), with real
coefficients, is positive semidefinite for all \(x\ge0\) if and only if
it can be written in the form
\[
  M_{ab}(x)
  =
  \sum_{i,j} Y^{ij}_{ab}\,q_i(x)q_j(x)
  +
  x\sum_{i,j} Z^{ij}_{ab}\,q_i(x)q_j(x),
\]
with \(Y,Z\) positive semidefinite matrices.  This converts the bootstrap
problem into a semidefinite program.

For the 3D Ising CFT, the exclusion plots in
\((\Delta_\sigma,\Delta_\epsilon)\) show a kink near
\[
  \Delta_\sigma\simeq0.518,\qquad
  \Delta_\epsilon\simeq1.412,
\]
and a lower bound on \(C_T\) is minimized near the same point.  This is
the numerical signature that the 3D Ising CFT sits at the corner of the
allowed region.

The same linear-functional method bounds OPE coefficients.  Writing the
identical-scalar crossing equation as
\[
  \sum_{\Delta,\ell}\lambda_{\Delta,\ell}^2
  H_{\Delta,\ell}=0,
  \qquad
  H_{\Delta,\ell}
  =
  v^{\Delta_\phi}G_{\Delta,\ell}(u,v)
  -u^{\Delta_\phi}G_{\Delta,\ell}(v,u),
\]
choose a functional \(\alpha\) such that
\[
  \alpha(H_{\Delta,\ell})\ge0\quad ((\Delta,\ell)\in I'),
  \qquad
  \alpha(H_{\Delta_*,\ell_*})=1.
\]
Then minimizing \(-\alpha(H_{\mathbf1})\) gives an upper bound on
\(\lambda_{\Delta_*,\ell_*}^2\).  For the stress tensor
\((\Delta_*,\ell_*)=(d,2)\),
\(\lambda_{\phi\phi T}^2\) is proportional to
\(\Delta_\phi^2/C_T\), up to convention-dependent constants, so this is
equivalently a lower bound on \(C_T\).  The notes quote a
three-dimensional Ising value near
\[
  \frac{C_T}{C_T^{\rm free\ scalar}}\simeq 0.946,
\]
and, from mixed-correlator bootstrap equations involving
\(\langle\sigma\sigma\sigma\sigma\rangle\),
\(\langle\sigma\sigma\epsilon\epsilon\rangle\), and
\(\langle\epsilon\epsilon\epsilon\epsilon\rangle\),
\[
  \Delta_\sigma=0.5181489(10),
  \qquad
  \Delta_\epsilon=1.412625(10).
\]

\begin{figure}[htbp]
\centering
\begin{tikzpicture}[scale=0.9,>=Stealth]
  % Source family: 253c pp. 103--112.
  \draw[->] (-5.1,-1.3) -- (-1.5,-1.3) node[right] {\(\Delta_\phi\)};
  \draw[->] (-5.1,-1.3) -- (-5.1,1.8) node[above] {\(\Delta_\epsilon\)};
  \draw[very thick,magenta] (-4.65,-0.75) .. controls (-4.2,0.25) and (-3.1,1.25) .. (-1.8,1.45);
  \draw[very thick,blue!70!black] (-4.55,-0.92) .. controls (-3.8,-0.15) and (-2.75,0.55) .. (-1.8,0.75);
  \draw[green!45!black,thick,pattern=north east lines,pattern color=green!35] (-4.9,-1.25) rectangle (-1.75,0.65);
  \node[magenta] at (-2.55,1.55) {excluded};
  \node[blue!70!black] at (-2.35,0.45) {allowed};
  \node[orange!85!black] at (-3.45,-0.55) {3D Ising kink};

  \draw[qft arrow] (-0.7,0.1) -- (0.55,0.1);

  \draw[->] (1.1,-1.3) -- (4.7,-1.3) node[right] {\(\Delta_\phi\)};
  \draw[->] (1.1,-1.3) -- (1.1,1.8) node[above] {\(C_T/C_T^{\rm free}\)};
  \draw[very thick,magenta] (1.35,1.35) .. controls (2.0,0.72) and (2.55,0.38) .. (3.2,0.28)
    .. controls (3.9,0.18) and (4.25,0.22) .. (4.45,0.25);
  \node[blue!70!black] at (3.15,0.02) {3D Ising CFT};
  \draw[blue!70!black,fill=blue!70!black] (3.0,0.32) circle (1.8pt);
  \node[qft note,orange!85!black] at (2.85,1.55)
    {linear functionals bound spectra and OPE data};
\end{tikzpicture}
\caption{Bootstrap exclusion curves and the Ising kink.}
\end{figure}

\section{A Survey of Conformal Field Theories}

The notes pause here to place the bootstrap examples in a broader
landscape of unitary, Poincar\'e-invariant local CFTs with a stress
tensor.  In \(d=1\) there is no ordinary notion of spacetime locality,
while in \(d\ge7\) interacting CFTs are often conjectured not to exist,
although this is not a theorem.

Free scalars and free fermions give basic examples in any dimension.
Free Maxwell theory is conformal in \(d=4\); in other dimensions it is
scale invariant but not a conformal field theory in the same sense.

In \(d=2\), the infinite-dimensional Virasoro symmetry leads to many
special families:
\begin{itemize}
  \item rational CFTs, often constructed from Wess--Zumino--Witten
  models and cosets;
  \item conformal nonlinear sigma models, with known interacting
  non-free examples relying on supersymmetry;
  \item noncompact CFTs with continuous spectra, including Liouville
  theory, Toda theory, the \(H_3^+\) model, and the cigar coset;
  \item orbifolds and infrared fixed points of RG flows, such as
  coupled Potts models.
\end{itemize}

In \(d=3\), important examples include Landau--Ginzburg fixed points
such as the critical \(O(N)\) model, fermionic fixed points such as the
Gross--Neveu model, Chern--Simons--matter theories, and massless gauge
theories such as QED\(_3\) with scalar or fermionic matter.  Holographic
duals of \(\mathrm{AdS}_4\) string vacua provide further examples, not
all of which have an identified field-theory description.

In \(d=4\), Banks--Zaks theories arise as weakly coupled fixed points of
QCD with \(N_f\) massless quarks in the conformal window,
\[
  \frac{34N_c^3}{13N_c^2-3}
  \lesssim N_f
  \lesssim
  \frac{11}{2}N_c,
\]
as estimated from the two-loop beta function.  Four dimensions also
contain a large class of supersymmetric conformal gauge theories,
including \(\mathcal N=1,2,4\) examples.

In \(d=5\), \(\mathcal N=1\) superconformal theories can have relevant
deformations flowing to five-dimensional supersymmetric gauge theories,
and possible nonsupersymmetric examples related to \(SU(2)\) gauge
theory have also been discussed.  In \(d=6\), \((2,0)\) and \((1,0)\)
superconformal theories are inferred from string-theoretic constructions
and supported by bootstrap constraints.

\section{Two-Dimensional CFT and Virasoro Algebra}

In two dimensions, introduce lightcone or complex coordinates
\[
  z=x^1+\ii x^2,\qquad \bar z=x^1-\ii x^2.
\]
The stress tensor decomposes into holomorphic and antiholomorphic pieces
\[
  T(z)=T_{zz}(z),\qquad \widetilde T(\bar z)=T_{\bar z\bar z}(\bar z),
\]
and the conformal transformations become infinite-dimensional.
The conformal Killing equations imply
\[
  z\mapsto f(z),\qquad
  \bar z\mapsto \bar f(\bar z),
\]
or infinitesimally
\[
  \epsilon(z)=\sum_n \epsilon_n z^{n+1},\qquad
  \bar\epsilon(\bar z)=\sum_n\bar\epsilon_n \bar z^{n+1}.
\]
The associated modes are
\[
  L_n=\oint_C\frac{\dd z}{2\pi\ii}\,z^{n+1}T(z),
  \qquad
  \widetilde L_n=\oint_C\frac{\dd\bar z}{2\pi\ii}\,
  \bar z^{n+1}\widetilde T(\bar z).
\]
For a primary field of weights \((h,\bar h)\),
\[
  \Delta=h+\bar h,\qquad s=h-\bar h,
\]
the stress-tensor OPE is
\[
  T(z)\phi(0)
  \sim
  \frac{h\,\phi(0)}{z^2}
  +\frac{\partial\phi(0)}{z},
  \qquad
  \widetilde T(\bar z)\phi(0)
  \sim
  \frac{\bar h\,\phi(0)}{\bar z^2}
  +\frac{\bar\partial\phi(0)}{\bar z}.
\]

\begin{figure}[htbp]
\centering
\begin{tikzpicture}[scale=0.9,>=Stealth]
  % Source family: 253c pp. 116--123.
  \draw[->] (-5.0,0) -- (-1.6,0) node[right] {\(\Re z\)};
  \draw[->] (-3.35,-1.55) -- (-3.35,1.55) node[above] {\(\Im z\)};
  \draw[qft contour] (-3.35,0) circle (0.95);
  \node[qft insertion] at (-3.35,0) {};
  \node[violet] at (-2.75,0.2) {\(\mathcal O(0)\)};
  \node[violet] at (-3.9,1.2) {\(C\)};
  \node at (-3.35,-2.0) {\(L_n=\oint_C \frac{\dd z}{2\pi i}z^{n+1}T(z)\)};

  \draw[qft arrow] (-0.9,0) -- (0.25,0);

  \draw[qft contour] (2.0,0) circle (1.05);
  \draw[qft contour,orange!85!black] (2.0,0) circle (0.65);
  \node[qft insertion] at (2.0,0) {};
  \node[violet] at (2.55,0.2) {\(\phi(0)\)};
  \node[violet] at (1.0,1.05) {\(C_1\)};
  \node[orange!85!black] at (2.95,0.72) {\(C_2\)};
  \draw[qft flow] (3.55,0.7) .. controls (4.15,0.25) and (4.15,-0.25) .. (3.55,-0.7);
  \node[qft note,blue!70!black] at (2.05,-1.65)
    {shrinking and expanding contours gives commutators};
\end{tikzpicture}
\caption{Virasoro modes as contour integrals of the stress tensor.}
\end{figure}

The OPE of the stress tensor with itself is
\[
  T(z)T(0)
  \sim
  \frac{c/2}{z^4}+\frac{2T(0)}{z^2}+\frac{\partial T(0)}{z},
\]
which is equivalent to the Virasoro algebra
\[
  [L_n,L_m]=(n-m)L_{n+m}
  +\frac{c}{12}(n^3-n)\delta_{n,-m}.
\]
Similarly,
\[
  [\widetilde L_n,\widetilde L_m]
  =
  (n-m)\widetilde L_{n+m}
  +\frac{\widetilde c}{12}(n^3-n)\delta_{n,-m},
  \qquad
  [L_n,\widetilde L_m]=0.
\]
The central term comes from the singular part of the \(T(z)T(w)\) OPE;
equivalently, it is the residue obtained when the \(L_n\) contour is
moved across the \(T(w)\) insertion.
The contour calculation is
\[
  [L_n,L_m]\phi(0)
  =
  \oint_{C_1}\frac{\dd z}{2\pi\ii}z^{n+1}T(z)
  \oint_{C_2}\frac{\dd w}{2\pi\ii}w^{m+1}T(w)\phi(0)
  -(n\leftrightarrow m),
\]
where \(C_1\) is deformed across \(C_2\).  After inserting the
\(T(z)T(w)\) OPE, the residue gives the Virasoro commutator above.
For a primary,
\[
  L_0\ket h=h\ket h,\qquad
  L_n\ket h=0\quad (n>0),
\]
and descendants are generated by \(L_{-n}\)'s; the same statements hold
for the antiholomorphic generators.

Under a holomorphic change of coordinates \(z\mapsto w(z)\), the stress
tensor transforms with a Schwarzian derivative:
\[
  T'(w)
  =
  \left(\frac{\dd z}{\dd w}\right)^2T(z)
  +\frac c{12}\{z,w\},
  \qquad
  \{z,w\}
  =
  \frac{z'''(w)}{z'(w)}
  -\frac32\left(\frac{z''(w)}{z'(w)}\right)^2 .
\]
For the plane-cylinder map \(z=\ee^{-\ii w}\),
\[
  T_{\rm cyl}(w)=-z^2T(z)+\frac c{24}.
\]
Thus the Hamiltonian on the spatial circle is shifted by the Casimir
energy:
\[
  H=L_0-\frac c{24}+\widetilde L_0-\frac{\widetilde c}{24},
  \qquad
  P=L_0-\frac c{24}-(\widetilde L_0-\frac{\widetilde c}{24}).
\]

\begin{figure}[htbp]
\centering
\begin{tikzpicture}[scale=0.88,>=Stealth]
  % Source family: 253c pp. 126--132.
  \draw[thick] (-5.1,-1.25) rectangle (-3.0,1.25);
  \node[qft insertion] at (-4.05,0) {};
  \node[violet] at (-4.05,0.35) {\(\phi(0)\)};
  \draw[qft arrow] (-2.35,0) -- (-1.1,0) node[midway,above] {\(z=\ee^{-iw}\)};
  \draw[thick] (0.1,1.25) ellipse (0.7 and 0.18);
  \draw[thick] (-0.6,1.25) -- (-0.6,-1.25);
  \draw[thick] (0.8,1.25) -- (0.8,-1.25);
  \draw[thick] (0.1,-1.25) ellipse (0.7 and 0.18);
  \draw[qft surface] (-0.6,0.2) arc[start angle=180,end angle=360,x radius=0.7,y radius=0.18];
  \draw[qft surface,dashed] (-0.6,0.2) arc[start angle=180,end angle=0,x radius=0.7,y radius=0.18];
  \node[violet] at (0.1,-0.35) {\(\ket{\phi}\)};
  \node[qft note,blue!70!black] at (0.15,-1.75)
    {plane--cylinder map shifts \(L_0\) by \(c/24\)};

  \draw[->] (2.0,-1.25) -- (5.2,-1.25) node[right] {\(h\)};
  \draw[->] (2.0,-1.25) -- (2.0,1.55);
  \draw[orange!85!black,thick] (2.45,-0.75) -- (3.35,0.1) -- (4.4,0.9);
  \draw[magenta,thick] (2.45,-0.25) -- (3.35,0.45) -- (4.4,1.15);
  \draw[green!45!black,thick,dashed] (3.0,-1.25) -- (3.0,1.35);
  \node[green!45!black] at (3.0,-1.55) {\(c=1\)};
  \node[orange!85!black] at (4.0,0.45) {unitarity};
  \node[magenta] at (4.05,1.35) {minimal models};
\end{tikzpicture}
\caption{Plane-cylinder map and the two-dimensional unitarity plot.}
\end{figure}

Null descendants produce differential equations for correlation
functions.  For example, the insertion of a degenerate field yields a
BPZ equation for a four-point function on the sphere.
At level two, a null state has the form
\[
  \left(L_{-2}-\frac{3}{2(2h+1)}L_{-1}^2\right)\ket h=0,
\]
with \(h,c\) related by the vanishing of the descendant norm.  Inserting
this field in a correlator and moving the \(T\)-contour onto the other
operators gives a second-order differential equation.  For a four-point
function
\[
  \left\langle
    \mathcal O_1(z_1,\bar z_1)\cdots
    \mathcal O_4(z_4,\bar z_4)
  \right\rangle
  =
  |z_{12}|^{-2\Delta_1}|z_{34}|^{-2\Delta_3} f(z,\bar z),
\]
the contour computation reduces the null-state relation to a
differential equation for \(f(z,\bar z)\).  The two independent local
solutions are exchanged by analytic continuation, and crossing fixes the
allowed linear combination.
Equivalently, for a degenerate insertion at \(z\),
\[
  \left[
    \frac{3}{2(2h+1)}\partial_z^2
    -
    \sum_i\left(
      \frac{h_i}{(z_i-z)^2}
      -\frac{1}{z_i-z}\partial_{z_i}
    \right)
  \right]
  \left\langle\phi_h(z)\prod_i\mathcal O_i(z_i)\right\rangle=0.
\]
At low levels the Gram matrices make the unitarity constraints explicit.
For the level-two descendants \(L_{-2}\ket h\) and \(L_{-1}^2\ket h\),
\[
  M^{(2)}
  =
  \begin{pmatrix}
    4h+\frac c2 & 6h\\
    6h & 4h(2h+1)
  \end{pmatrix},
\]
and positivity of such matrices at every level leads to the two
branches quoted next.

Unitarity in two dimensions gives either
\[
  c\ge1,\qquad h,\bar h\ge0,
\]
or the minimal-model series
\[
  c=1-\frac{6}{m(m+1)},\qquad m=3,4,5,\ldots,
\]
with
\[
  h_{r,s}
  =
  \frac{\bigl((m+1)r-ms\bigr)^2-1}{4m(m+1)},
  \qquad
  1\le r\le m-1,\quad 1\le s\le r.
\]
This triangular range chooses one representative from the Kac-table
identification.

\begin{figure}[htbp]
\centering
\begin{tikzpicture}[scale=0.9,>=Stealth]
  % Source family: 253c pp. 130--135.
  \draw[thick] (-4.2,0) ellipse (1.75 and 1.05);
  \foreach \x/\y/\lab in {-5.05/0.35/{z_1},-3.55/0.45/{z_2},-4.65/-0.55/{z_3},-3.35/-0.52/{z_4}} {
    \node[qft insertion] at (\x,\y) {};
    \node[above right,violet] at (\x,\y) {\(\lab\)};
  }
  \draw[qft contour] (-5.05,0.35) circle (0.32);
  \draw[qft contour] (-3.55,0.45) circle (0.32);
  \draw[qft contour] (-4.65,-0.55) circle (0.32);
  \draw[qft contour] (-3.35,-0.52) circle (0.32);
  \node[qft note,blue!70!black] at (-4.2,-1.55)
    {contours around insertions become BPZ differential operators};

  \node at (-0.9,0) {\(\Longrightarrow\)};

  \node at (2.0,0.35) {\(\displaystyle
    \left[\partial_z^2+
    \sum_i\left(\frac{h_i}{(z-z_i)^2}
    +\frac{\partial_{z_i}}{z-z_i}\right)\right]f(z_i)=0\)};
  \draw[thick,blue!70!black] (0.05,-0.55) -- (4.15,-0.55);
  \node[orange!85!black] at (2.05,-1.05) {null-state decoupling};
\end{tikzpicture}
\caption{BPZ equations from contour deformation and null states.}
\end{figure}

The first nontrivial unitary minimal model is \(m=3\), for which
\(c=1/2\).  Its allowed holomorphic weights are
\[
  h_{1,1}=0,\qquad h_{2,1}=\frac12,\qquad h_{2,2}=\frac1{16}.
\]
Together with the antiholomorphic sector, these give the primaries of
the two-dimensional Ising CFT:
\[
  \mathbf1,\qquad
  \epsilon\ \text{with}\ (h,\bar h)=\left(\frac12,\frac12\right),
  \qquad
  \sigma\ \text{with}\ (h,\bar h)=\left(\frac1{16},\frac1{16}\right).
\]
The spin field \(\sigma\) has a level-two null descendant,
\[
  \left(L_{-2}-\frac43L_{-1}^2\right)\ket\sigma=0,
  \qquad
  L_{-2}\cdot\sigma=\frac43\,\partial^2\sigma ,
\]
and similarly in the antiholomorphic sector.  Applying the BPZ contour
argument to
\[
  \langle\sigma(z_1,\bar z_1)\cdots\sigma(z_4,\bar z_4)\rangle
  =
  |z_{12}|^{-2\Delta_\sigma}|z_{34}|^{-2\Delta_\sigma} f(z,\bar z)
\]
gives
\[
  \left[
    \partial_z^2
    +\frac{2-5z}{4z(1-z)}\partial_z
    -\frac{3}{64(1-z)^2}
  \right]f(z,\bar z)=0,
\]
with an identical equation in \(\bar z\).  Two holomorphic solutions are
\[
  f_\pm(z)=(1-z)^{-1/8}\sqrt{1\pm\sqrt z}.
\]
Single-valuedness around \(z=0\) and \(z=1\), together with
normalization \(f(0,0)=1\), fixes
\[
  f(z,\bar z)
  =
  \frac{|1+\sqrt z|+|1-\sqrt z|}
       {2|1-z|^{1/4}}.
\]
Crossing under \(z\mapsto1-z\) follows from the elementary identity
\[
  \bigl(|1+w|+|1-w|\bigr)^2
  =
  2+2|w|^2+2|1-w^2|,
\]
applied to \(w=\sqrt z\) and \(w=\sqrt{1-z}\).  Thus the level-two null
state determines a crossing-symmetric four-point function.

Comparing with the Virasoro block expansion,
\[
  f(z,\bar z)
  =
  |\mathcal F_0(z)|^2
  +
  C_{\sigma\sigma\epsilon}^2\,|\mathcal F_{1/2}(z)|^2,
\]
one may take
\[
  \mathcal F_0(z)=\frac{f_+(z)+f_-(z)}2,
  \qquad
  \mathcal F_{1/2}(z)=f_+(z)-f_-(z),
\]
so that
\[
  C_{\sigma\sigma\epsilon}=\frac12.
\]
The other scalar structure constants quoted in the notes are
\[
  C_{\epsilon\epsilon\sigma}=0,\qquad
  C_{\epsilon\epsilon\epsilon}=0,
\]
the first by \(\mathbb Z_2\) symmetry and the second by the additional
selection rule of the Ising model.

\section{Modular Invariance and Weyl Anomaly}

For a two-dimensional CFT on a circle, the cylinder Hamiltonian is
\[
  H=L_0-\frac c{24}+\widetilde L_0-\frac{\widetilde c}{24}.
\]
The torus partition function is
\[
  Z(\tau)=\operatorname{Tr}\,
  \exp\left(2\pi\ii\tau\left(L_0-\frac c{24}\right)
  -2\pi\ii\bar\tau\left(\widetilde L_0-\frac{\widetilde c}{24}\right)\right).
\]
Euclidean locality allows one to interchange the two cycles of the
torus.  Therefore
\[
  Z(\tau)=Z(-1/\tau),
  \qquad
  Z(\tau)=Z(\tau+1),
\]
and hence
\[
  Z(\tau)=Z\left(\frac{a\tau+b}{c\tau+d}\right),
  \qquad
  a,b,c,d\in\mathbb Z,\quad ad-bc=1.
\]
This modular invariance is a highly nontrivial constraint on the CFT
spectrum because it relates low-temperature and high-temperature
expansions of the same partition function.

\begin{figure}[htbp]
\centering
\begin{tikzpicture}[scale=0.9,>=Stealth]
  % Source family: 253c pp. 141--144.
  \draw[thick] (-5.35,1.15) ellipse (0.85 and 0.22);
  \draw[thick] (-6.2,1.15) -- (-6.2,-1.15);
  \draw[thick] (-4.5,1.15) -- (-4.5,-1.15);
  \draw[thick] (-5.35,-1.15) ellipse (0.85 and 0.22);
  \draw[qft arrow] (-5.95,-1.35) -- (-5.95,1.35) node[midway,left] {\(\beta\)};
  \draw[qft arrow] (-6.25,-1.45) -- (-4.45,-1.45) node[midway,below] {\(2\pi\)};
  \node at (-5.35,-1.9) {thermal cylinder};

  \draw[qft arrow] (-3.6,0) -- (-2.45,0);

  \draw[thick] (-1.2,-0.8) -- (0.8,-0.8) -- (1.35,0.85) -- (-0.65,0.85) -- cycle;
  \draw[blue!70!black,thick] (-1.2,-0.8) -- (0.8,-0.8);
  \draw[green!45!black,thick] (-1.2,-0.8) -- (-0.65,0.85);
  \node[blue!70!black] at (-0.2,-1.12) {\(2\pi\)};
  \node[green!45!black] at (-1.25,0.12) {\(2\pi\tau\)};
  \node at (0.05,-1.9) {torus lattice};

  \draw[qft arrow] (2.0,0) -- (3.1,0);
  \draw[thick] (4.25,-0.75) -- (5.8,-0.75) -- (5.3,0.75) -- (3.75,0.75) -- cycle;
  \draw[thick,orange!85!black] (4.25,-0.75) -- (5.3,0.75);
  \draw[thick,violet] (3.75,0.75) -- (5.8,-0.75);
  \node[orange!85!black] at (5.15,1.05) {\(S:\tau\mapsto -1/\tau\)};
  \node[violet] at (4.45,-1.1) {\(T:\tau\mapsto \tau+1\)};
\end{tikzpicture}
\caption{Thermal cylinder, torus lattice, and modular transformations.}
\end{figure}

The Weyl anomaly in two dimensions is encoded in the response of the
effective action to a Weyl rescaling,
\[
  \delta_\omega W[g]=-\frac c{24\pi}\int \dd^2x\sqrt g\,\omega R.
\]
Equivalently,
\[
  T^\mu{}_\mu=-\frac c{12}R
\]
in the normalization of the notes.
The stress tensor insertion follows from metric variation:
\[
  \delta S
  =
  \frac12\int\dd^2x\sqrt g\,\delta g^{\mu\nu}T_{\mu\nu}.
\]
Locally, a two-dimensional metric deformation decomposes into a
diffeomorphism plus a Weyl transformation,
\[
  \delta g_{\mu\nu}
  =
  \nabla_\mu\epsilon_\nu+\nabla_\nu\epsilon_\mu
  +2\delta\omega\,g_{\mu\nu}.
\]
Classically, the CFT is invariant under both pieces.  Quantum
mechanically, contact terms in the overlapping stress-tensor insertions
produce an anomaly.  In flat coordinates,
\[
  T(z)T(0)
  \sim
  \frac{c/2}{z^4}
  +\frac{2T(0)}{z^2}
  +\frac{\partial T(0)}{z},
\]
and the polar terms should be interpreted distributionally; for example
\[
  \bar\partial\frac1z=2\pi\delta^{(2)}(z)
\]
in the convention of the notes.  Requiring consistency of the contact
terms in the conservation equation gives \(c=\widetilde c\), the
absence of a gravitational anomaly.

For a finite Weyl transformation \(g'_{\mu\nu}=\ee^{2\omega}g_{\mu\nu}\),
the partition function transforms as
\[
  Z[g']=\ee^{-S_{\rm W}[\omega;g]}Z[g],
\]
where the local Wess-Zumino or dilaton action is
\[
  S_{\rm W}[\omega;g]
  =
  \frac c{24\pi}
  \int\dd^2x\sqrt g\,
  \left(g^{\mu\nu}\partial_\mu\omega\partial_\nu\omega+R\omega\right).
\]
With this understood, correlation functions on any Riemann surface can
be determined from conformal data on the plane.

\begin{figure}[htbp]
\centering
\begin{tikzpicture}[scale=0.88,>=Stealth]
  % Source family: 253c pp. 151--153.
  \draw[thick] (-5.4,0) .. controls (-4.9,0.85) and (-4.0,0.9) .. (-3.45,0.25)
    .. controls (-2.85,-0.45) and (-2.0,-0.4) .. (-1.55,0.2)
    .. controls (-1.2,-0.75) and (-2.25,-1.15) .. (-3.15,-0.85)
    .. controls (-4.15,-0.55) and (-4.9,-1.05) .. (-5.4,0);
  \node[qft insertion] at (-4.85,0.25) {};
  \node[qft insertion] at (-2.15,-0.2) {};
  \node[violet] at (-5.1,0.65) {\(z_i\)};
  \node[violet] at (-1.9,0.2) {\(z_j\)};
  \draw[qft arrow] (-0.9,0) -- (0.25,0) node[midway,above] {Weyl};
  \draw[thick] (1.2,-1.05) rectangle (4.5,1.05);
  \foreach \x/\y/\lab in {1.85/0.45/{z_i},3.8/-0.35/{z_j}} {
    \node[qft insertion] at (\x,\y) {};
    \node[above right,violet] at (\x,\y) {\(\lab\)};
  }
  \draw[qft contour] (2.0,0.0) circle (0.42);
  \draw[qft contour] (3.65,0.0) circle (0.42);
  \node[qft note,blue!70!black] at (1.45,-1.5)
    {curved background correlators are determined by flat-space data plus anomaly};
\end{tikzpicture}
\caption{Weyl map from a curved surface to local flat coordinates.}
\end{figure}

Modular invariance can also be visualized by cutting a Riemann surface
into pairs of pants.  Different decompositions of the same surface give
the same answer and are related to crossing and modular transformations.
Any two pair-of-pants decompositions are related by a sequence of basic
crossing moves and modular moves.  Thus the four-point crossing
equation, together with modular invariance of the torus one-point
channel, is the local form of full consistency of the CFT on arbitrary
two-dimensional surfaces.

\begin{figure}[htbp]
\centering
\begin{tikzpicture}[scale=0.9,>=Stealth]
  % Source family: 253c pp. 152--153.
  \draw[thick] (-5.55,0.2) .. controls (-5.05,1.1) and (-3.9,1.1) .. (-3.45,0.25)
    .. controls (-3.1,-0.42) and (-3.75,-0.92) .. (-4.55,-0.7)
    .. controls (-5.25,-0.5) and (-5.95,-0.6) .. (-5.55,0.2);
  \draw[thick] (-3.45,0.25) .. controls (-2.95,1.0) and (-1.95,0.95) .. (-1.55,0.1)
    .. controls (-1.15,-0.65) and (-2.05,-0.95) .. (-2.75,-0.65)
    .. controls (-3.25,-0.42) and (-3.78,-0.32) .. (-3.45,0.25);
  \draw[qft contour] (-3.35,0.0) circle (0.38);
  \node at (-3.5,-1.25) {pair-of-pants channel 1};

  \node at (-0.05,0) {\(=\)};

  \draw[thick] (1.05,0.25) .. controls (1.7,1.25) and (3.35,1.15) .. (3.7,0.15)
    .. controls (4.1,-0.85) and (2.75,-1.05) .. (2.15,-0.55)
    .. controls (1.6,-0.05) and (0.45,-0.65) .. (1.05,0.25);
  \draw[thick] (4.05,0.25) .. controls (4.65,1.05) and (5.75,0.8) .. (5.9,-0.05)
    .. controls (6.0,-0.85) and (4.75,-0.95) .. (4.25,-0.45)
    .. controls (3.82,-0.05) and (3.6,-0.25) .. (4.05,0.25);
  \draw[qft contour,orange!85!black] (3.75,0.05) circle (0.38);
  \node at (3.6,-1.25) {pair-of-pants channel 2};
\end{tikzpicture}
\caption{Different pair-of-pants decompositions give the same surface amplitude.}
\end{figure}

\section{Supersymmetry and \texorpdfstring{\(N=4\)}{N=4} Yang--Mills}

Consider first a four-dimensional gauge theory with massless matter,
with action schematically
\[
  S=\int\dd^4x\,
  \left[
    -\frac1{2g^2}\tr F_{\mu\nu}F^{\mu\nu}
    -\bar\psi\gamma^\mu D_\mu\psi
    -(D_\mu\phi)^\dagger D^\mu\phi+\cdots
  \right],
\]
where
\[
  F_{\mu\nu}=\partial_\mu A_\nu-\partial_\nu A_\mu-\ii[A_\mu,A_\nu],
  \qquad
  A_\mu=A_\mu^a t^a,
  \qquad
  \tr(t^at^b)=\frac12\delta^{ab}.
\]
The renormalized coupling is defined through the background-field
one-particle-irreducible effective action,
\[
  \Gamma[A]
  =
  -\frac1{2g^2(\mu)}
  \int\dd^4x\,\tr F_{\mu\nu}F^{\mu\nu}
  +\cdots .
\]
In dimensional regularization and minimal subtraction,
\[
  \mu\frac{\dd g}{\dd\mu}=\beta(g),
  \qquad
  \beta(g)=-\frac{b}{(4\pi)^2}g^3+O(g^5).
\]
For \(G=SU(N)\), with \(N_f\) Weyl fermions and \(N_b\) real scalars in
the adjoint representation,
\[
  b=\frac N3(11-2N_f-\tfrac12N_b).
\]
Thus the field content \(N_f=4\), \(N_b=6\) makes the one-loop
coefficient vanish.  This is the field content of \(N=4\) super
Yang--Mills, where the vanishing extends to the exact beta function.

The \(N=4\) super Yang--Mills theory contains
\[
  A_\mu,\qquad
  \lambda^I_\alpha\quad(I=1,\ldots,4),\qquad
  \phi^i\quad(i=1,\ldots,6),
\]
all in the adjoint representation of \(SU(N)\).  The global
R-symmetry is
\[
  SO(6)\simeq SU(4),
\]
under which the fermions transform as the fundamental spinor of
\(SU(4)\), while the six scalars transform as the vector of \(SO(6)\).
The Lagrangian is schematically
\[
\begin{aligned}
  \mathcal L
  =\tr\bigg[
    &-\frac1{2g^2}F_{\mu\nu}F^{\mu\nu}
    -\frac1{g^2}D_\mu\phi^iD^\mu\phi^i
    -\frac1{g^2}\bar\lambda_I\gamma^\mu D_\mu\lambda^I  \\
    &+\frac1{g^2}[\phi^i,\phi^j]^2
    +\frac1{g^2}\lambda\lambda\phi
  \bigg],
\end{aligned}
\]
with the relative coefficients fixed by supersymmetry.  Although written
schematically here, the important point in the notes is that the Yukawa
couplings and scalar potential are tied to the gauge coupling, so their
RG running is tied to \(\beta_g\).

In supersymmetric theories, supercharges relate bosonic and fermionic
operators.  The state/operator map again turns local operators into
states on \(S^{d-1}\), and superconformal primaries are annihilated by
the special supercharges \(S\).
The Poincare supercharges obey, in spinor notation,
\[
  \{Q^I_\alpha,\bar Q_{\dot\beta J}\}
  =
  2\delta^I{}_J\,\sigma^\mu_{\alpha\dot\beta}P_\mu,
  \qquad
  [P_\mu,Q]=0.
\]
Because
\[
  [P_\mu,K_\nu]=2\ii(\eta_{\mu\nu}D-J_{\mu\nu}),
\]
the special conformal generators cannot commute with \(Q\).  Their
commutators define the special supercharges \(S,\bar S\):
\[
  [K_\mu,Q^I_\alpha]\sim(\sigma_\mu)_{\alpha\dot\alpha}\bar S^{\dot\alpha I},
  \qquad
  [K_\mu,\bar Q_{\dot\alpha I}]\sim
  (\bar\sigma_\mu)^{\dot\alpha\alpha}S_{\alpha I}.
\]
The mixed anticommutator has the schematic form
\[
  \{Q,S\}\sim D+J+R,
\]
where \(R\) denotes the \(SU(4)\) R-symmetry generators.  After Wick
rotation and the state/operator map, \(Q\) and \(S\) are Hermitian
conjugates on the Hilbert space on \(S^3\), just as \(P_\mu\) and
\(K_\mu\) are conjugates.  Positivity of norms of \(Q\)-descendants
therefore gives the BPS bounds.

\begin{figure}[htbp]
\centering
\begin{tikzpicture}[scale=0.9,>=Stealth]
  % Source family: 253c pp. 157--163.
  \draw[thick] (-5.1,-1.25) rectangle (-3.0,1.25);
  \draw[qft surface] (-4.05,0) circle (0.78);
  \node[qft insertion] at (-4.05,0) {};
  \node[violet] at (-4.45,0.38) {\(\mathcal O(0)\)};
  \draw[qft arrow] (-2.35,0) -- (-1.1,0) node[midway,above] {state/operator};
  \draw[thick] (0.1,1.25) ellipse (0.75 and 0.2);
  \draw[thick] (-0.65,1.25) -- (-0.65,-1.25);
  \draw[thick] (0.85,1.25) -- (0.85,-1.25);
  \draw[thick] (0.1,-1.25) ellipse (0.75 and 0.2);
  \node[violet] at (0.1,-0.25) {\(\ket{\mathcal O}\)};
  \draw[qft flow] (0.1,0.1) -- (0.1,0.9) node[midway,right] {\(Q\)};
  \draw[qft flow,orange!85!black] (0.1,0.1) -- (0.75,-0.65) node[midway,right] {\(S\)};

  \draw[qft arrow] (1.2,0) -- (1.9,0);
  \node[anchor=west,align=left] at (2.15,0.5)
    {\(\{Q,S\}\sim D+J+R\)\\[0.25em]
     \(\Delta\ge\) linear function of spins and \(R\)-charges};
  \node[qft note,blue!70!black] at (3.95,-0.85)
    {BPS shortening occurs when a descendant has zero norm};
\end{tikzpicture}
\caption{Superconformal descendants and BPS shortening.}
\end{figure}

For \(N=4\) super Yang--Mills, the fields transform under an \(SU(4)\)
R-symmetry.  Highest weights can be described by Dynkin labels
\([r_1,r_2,r_3]\), and BPS bounds constrain the scaling dimension.
Choose Cartan generators \(H_1,H_2,H_3\) of \(SU(4)\).  In any
representation, one can diagonalize them simultaneously; the resulting
triples are weights.  The remaining generators shift weights by roots.
For \(SU(4)\), every irreducible representation is labeled by a highest
weight, equivalently by Dynkin labels
\[
  [r_1,r_2,r_3],
  \qquad r_i\in\mathbb Z_{\ge0}.
\]
The fundamental representations have highest weights corresponding to
the labels
\[
  \mathbf4=[1,0,0],
  \qquad
  \overline{\mathbf4}=[0,0,1],
  \qquad
  \mathbf6=[0,1,0],
\]
and the adjoint is
\[
  \mathbf{15}=[1,0,1].
\]

For a superconformal primary labeled by Lorentz spins
\((j,\bar j)\) and \(SU(4)\) Dynkin labels \([r_1,r_2,r_3]\), the notes
quote the schematic pair of unitarity/BPS alternatives
\[
  \Delta\ge
  2+j+\frac{3r_1+2r_2+r_3}{2}
  \quad(j>0),
  \qquad
  \Delta=
  \frac{3r_1+2r_2+r_3}{2}
  \quad(j=0),
\]
and similarly
\[
  \Delta\ge
  2+\bar j+\frac{r_1+2r_2+3r_3}{2}
  \quad(\bar j>0),
  \qquad
  \Delta=
  \frac{r_1+2r_2+3r_3}{2}
  \quad(\bar j=0).
\]
A representation is shortened when some \(Q,\bar Q\)-descendants are
null; this occurs when a BPS bound is saturated.

An important example is
\[
  \mathcal O_{ij}
  =
  \tr(\phi_i\phi_j)-\frac16\delta_{ij}\tr(\phi_k\phi_k),
\]
which transforms in the \(\mathbf{20}'=[0,2,0]\) of \(SO(6)_R\).  This
operator is a superconformal primary with protected dimension
\[
  \Delta=2
\]
at all values of the Yang--Mills coupling.  By contrast, the singlet
\(\tr(\phi_k\phi_k)\) belongs to a long representation and has a
nontrivial anomalous dimension.

\begin{figure}[htbp]
\centering
\begin{tikzpicture}[scale=0.9,>=Stealth]
  % Source family: 253c pp. 164--170.
  \foreach \x/\lab in {-3/{r_1},0/{r_2},3/{r_3}} {
    \draw[thick] (\x,0) circle (0.42);
    \node at (\x,0) {\(\lab\)};
  }
  \draw[thick] (-2.58,0) -- (-0.42,0);
  \draw[thick] (0.42,0) -- (2.58,0);
  \node[orange!85!black] at (0,0.9) {\(SU(4)\) Dynkin labels};
  \node at (0,-0.95) {\(\Delta\ge 2+j+\frac{3r_1+2r_2+r_3}{2}\)};
  \draw[qft arrow] (4.0,0.45) -- (5.0,0.45);
  \node[right] at (5.05,0.45) {long multiplet};
  \draw[qft arrow,green!45!black] (4.0,-0.45) -- (5.0,-0.45);
  \node[right,green!45!black] at (5.05,-0.45) {BPS shortening};
\end{tikzpicture}
\caption{R-symmetry labels and schematic superconformal shortening.}
\end{figure}

\section{Large-\texorpdfstring{\(N\)}{N} Expansion and Spin Chains}

In the large-\(N\) limit of gauge theory, Feynman diagrams are drawn as
ribbon graphs.  Each closed index loop contributes a factor of \(N\), so
planar diagrams dominate the expansion.
For a typical \(N\times N\) matrix field \(\Phi\), with
\[
  \mathcal L=\tr\left((\partial\Phi)^2+V(\Phi)+\cdots\right)
\]
invariant under \(\Phi\mapsto U^\dagger\Phi U\), the double-line
propagator carries two color indices.  A fat graph with \(V\) vertices,
\(E\) edges, and \(F\) index loops contributes a power
\[
  N^{F}g^{2E-2V}
  =
  N^{\chi}\lambda^{E-V},
  \qquad
  \lambda=g^2N,
  \qquad
  \chi=V-E+F=2-2h.
\]
Thus, in the 't Hooft limit
\[
  N\to\infty,\qquad \lambda=g^2N\ \text{fixed},
\]
graphs of genus \(h\) are suppressed by \(N^{-2h}\), and the planar
graphs dominate.

In pure Yang--Mills the RG equation for the 't Hooft coupling has the
form
\[
  \frac{\dd\lambda}{\dd\log\mu}
  =
  -\frac{b}{24\pi^2}\lambda^2+O(\lambda^3)
\]
in the normalization of the notes, leading to dimensional transmutation
and a scale \(\Lambda_{\rm QCD}\).  One expects glueball masses and flux
string tensions of order this scale,
\[
  M_{\rm glueball}\sim\Lambda_{\rm QCD},
  \qquad
  T_{\rm flux}\sim\Lambda_{\rm QCD}^2.
\]
Large \(N\) makes glueballs and flux strings weakly interacting: a
three-glueball effective coupling is suppressed by \(1/N\).

\begin{figure}[htbp]
\centering
\begin{tikzpicture}[scale=0.9,>=Stealth]
  % Source family: 253c pp. 171--178.
  \begin{scope}[shift={(-4.65,0)}]
    \draw[qft feynman] (-1.0,0) .. controls (-0.5,0.8) and (0.5,0.8) .. (1.0,0);
    \draw[qft feynman] (-1.0,0) .. controls (-0.5,-0.8) and (0.5,-0.8) .. (1.0,0);
    \draw[qft feynman] (-0.35,0.58) -- (-0.35,-0.58);
    \draw[qft feynman] (0.35,0.58) -- (0.35,-0.58);
    \node at (0,-1.25) {ordinary graph};
  \end{scope}
  \node at (-2.2,0) {\(=\)};
  \begin{scope}[shift={(-0.2,0)}]
    \draw[thick] (-1.2,0) .. controls (-0.65,0.9) and (0.65,0.9) .. (1.2,0);
    \draw[thick] (-1.2,0) .. controls (-0.65,-0.9) and (0.65,-0.9) .. (1.2,0);
    \draw[thick] (-0.85,0.22) .. controls (-0.35,0.58) and (0.35,0.58) .. (0.85,0.22);
    \draw[thick] (-0.85,-0.22) .. controls (-0.35,-0.58) and (0.35,-0.58) .. (0.85,-0.22);
    \draw[qft contour,blue!70!black] (0,0) circle (0.45);
    \node at (0,-1.25) {fat graph};
  \end{scope}
  \draw[qft arrow] (1.85,0) -- (3.0,0);
  \begin{scope}[shift={(4.3,0)}]
    \draw[thick,fill=blue!8] (-1.0,-0.65) rectangle (1.0,0.65);
    \foreach \x in {-0.65,-0.2,0.25,0.7} {
      \draw[qft feynman] (\x,-0.65) -- (\x,0.65);
    }
    \node at (0,-1.25) {planar index sheet};
  \end{scope}
\end{tikzpicture}
\caption{The large-\(N\) expansion rewrites diagrams as ribbon graphs.}
\end{figure}

Wilson loops create flux strings.  In the planar limit, single-trace
operators behave like closed spin chains whose sites are elementary
fields.
The Wilson loop is
\[
  W_R(C)
  =
  \tr_R\,P\exp\left(\ii\oint_C A_\mu\,\dd x^\mu\right),
\]
and correlators of Wilson loops measure splitting and joining amplitudes
of flux strings.  This is the sense in which large-\(N\) Yang--Mills is
a theory of weakly interacting strings.

Back in \(N=4\) SYM, single-trace and multi-trace operators decouple at
leading order in the planar limit.  To study operators of finite
dimension as \(N\to\infty\), it suffices to focus on single traces.  In
the \(SU(2)\) sector, choose two complex scalars
\[
  Z,\qquad X,
\]
and consider operators
\[
  \tr(A_1A_2\cdots A_L),
  \qquad
  A_i\in\{Z,X\}.
\]
The operator \(\tr Z^L\) is BPS and has \(\Delta=L\), while a word with
some \(X\)'s is interpreted as a spin chain with impurities.

\begin{figure}[htbp]
\centering
\begin{tikzpicture}[scale=0.9,>=Stealth]
  % Source family: 253c pp. 176--185.
  \draw[qft contour] (-5.1,0.55) .. controls (-4.25,1.15) and (-3.1,1.1) .. (-2.55,0.45)
    .. controls (-3.0,-0.4) and (-4.25,-0.45) .. (-5.1,0.55);
  \draw[qft contour,orange!85!black] (-4.95,-0.95) .. controls (-4.2,-0.35) and (-3.05,-0.45) .. (-2.65,-1.15);
  \node[violet] at (-4.05,1.35)
    {\(W_R(C)=\tr_R P\exp\!\left(\ii\oint_C A_\mu\,\dd x^\mu\right)\)};
  \node[orange!85!black] at (-3.75,-1.45) {flux string};

  \draw[qft arrow] (-1.6,0) -- (-0.5,0);

  \draw[thick] (1.2,0) ellipse (1.25 and 0.5);
  \foreach \ang/\lab in {20/{Z},80/{Z},140/{X},200/{Z},260/{Z},320/{Z}} {
    \coordinate (p) at ({1.2+1.25*cos(\ang)},{0.5*sin(\ang)});
    \node[blue!70!black] at (p) {\(\lab\)};
  }
  \draw[qft contour,magenta] (1.2,0) ellipse (1.25 and 0.5);
  \node at (1.2,-1.2) {single trace as spin chain};

  \draw[qft arrow] (3.0,0) -- (4.05,0);

  \draw[thick] (5.15,0) ellipse (1.0 and 0.42);
  \foreach \ang/\lab in {35/{p_1},145/{p_2},250/{}} {
    \coordinate (p) at ({5.15+1.0*cos(\ang)},{0.42*sin(\ang)});
    \node[magenta] at (p) {\(\lab\)};
    \draw[magenta,thick] (p) circle (0.12);
  }
  \node at (5.15,-1.2) {magnons};
  \node[qft note,blue!70!black] at (2.95,1.45)
    {\(\Delta-J\) is the spin-chain Hamiltonian};
\end{tikzpicture}
\caption{Wilson loops, single-trace spin chains, and magnon excitations.}
\end{figure}

At one loop in the \(SU(2)\) sector, the dilatation operator takes the
form of a nearest-neighbor spin-chain Hamiltonian,
\[
  \Delta_1
  =
  \frac{1}{8\pi^2}
  \sum_{\ell=1}^L(1-P_{\ell,\ell+1}).
\]
The permutation \(P_{\ell,\ell+1}\) exchanges the fields at neighboring
sites.  The missing ``diagonal'' one-loop diagrams are fixed by the BPS
condition that \(\tr Z^L\) remain protected.  A single impurity at site
\(\ell\) defines a magnon basis state; momentum eigenstates are
\[
  \ket p\sim\sum_{\ell=1}^L
  \ee^{\ii p\ell}\tr(Z^{\ell-1}XZ^{L-\ell}),
  \qquad
  p=\frac{2\pi n}{L}.
\]
Acting with the one-loop Hamiltonian gives
\[
  \Delta_1\ket p
  =
  \frac1{2\pi^2}\sin^2\frac p2\,\ket p,
\]
so, for
\[
  \Delta=\Delta_0+\lambda\Delta_1+O(\lambda^2),
\]
the one-loop dispersion relation is
\[
  \Delta(p)=1+\frac{\lambda}{2\pi^2}\sin^2\frac p2+O(\lambda^2).
\]
Single-magnon states have momentum \(p=2\pi n/L\), and the perturbative
dispersion relation is corrected into the exact form
\[
  \Delta(p)
  =
  \sqrt{1+\frac{\lambda}{\pi^2}\sin^2\frac p2}.
\]
The exact expression is determined by an enhanced centrally extended
\(SU(2|2)\) symmetry of the spin chain.  The notes end with the
conjectural integrability statement: impurity scattering is governed by
a factorized exact S-matrix, which in principle determines the anomalous
dimensions on the planar single-trace spectrum.
